\section{Quantum double models}%
\label{sec:peps_quantum_double_models}

The quantum double model of $G$ has two PEPS representations
\cite[Appendix~A, ``The Toric Code and quantum double models'']{Cirac2021Matrix}.
After two facts on the labels of closures, this section constructs the dual
tensor, shows that it is $G$-isometric and that its network is the
superposition of the configurations obeying Gauss' law, and then treats the
primal tensor and the toric code.

\begin{lemma}[Pair conjugacy with a fixed first element]%
    \label{thm:peps_pair_conjugacy_centralizer}
    \lean{TNLean.PEPS.pairConjugacyClass_eq_iff_centralizer}
    \leanok
    \uses{def:peps_pair_conjugacy}
    In the fixed-first-element classification following
    \cite[Theorem~5.9]{Schuch2010PEPS}, fix $g\in G$. The pairs $(g,h)$ and $(g,k)$ have the same class
    exactly when $h=xkx^{-1}$ for some $x$ in the centralizer
    $Z_G(g)=\{x\in G:xg=gx\}$.
\end{lemma}

\begin{proof}\leanok
    The simultaneous-conjugacy equations are $g=xgx^{-1}$ and
    $h=xkx^{-1}$; the first is equivalent to $xg=gx$.
\end{proof}

\begin{lemma}[Closure labels for an abelian group]%
    \label{thm:peps_abelian_pair_classes}
    \lean{TNLean.PEPS.pairConjugacyClass_eq_iff_of_commGroup,
        TNLean.PEPS.pairConjugacyClassEquivOfCommGroup,
        TNLean.PEPS.commutingPairConjugacyClassEquivOfCommGroup,
        TNLean.PEPS.card_commutingPairConjugacyClass_of_commGroup}
    \leanok
    \uses{def:peps_pair_conjugacy}
    If $G$ is abelian, every pair-conjugacy class is a singleton and all
    pairs commute. A finite abelian group therefore has $|G|^2$
    commuting pair-conjugacy classes.
\end{lemma}

\begin{proof}\leanok
    Conjugation is trivial, so the class map identifies the classes with
    $G\times G$; counting pairs gives $|G|^2$.
\end{proof}

\begin{definition}[Dual quantum-double tensor]%
    \label{def:peps_quantum_double_dual}
    \lean{TNLean.PEPS.colorDiag, TNLean.PEPS.colorDiag_one,
        TNLean.PEPS.colorDiag_mul, TNLean.PEPS.quantumDoubleDualSpins,
        TNLean.PEPS.quantumDoubleDualTensor, TNLean.PEPS.colorShiftRep,
        TNLean.PEPS.colorShiftRep_apply,
        TNLean.PEPS.mem_invariants_colorShiftRep_iff}
    \leanok
    \uses{def:peps_g_injective}
    The virtual legs carry colors $h_1$ (left), $h_2$ (top), $h_3$ (right)
    and $h_4$ (down) in $G$, and the dual tensor is $1$ at the physical
    configuration
    $(h_2h_1^{-1},\,h_3h_2^{-1},\,h_3h_4^{-1},\,h_4h_1^{-1})$ of differences
    of adjacent colors and $0$ elsewhere \cite[Appendix~A, ``The Toric Code
        and quantum double models'']{Cirac2021Matrix}. The group acts on the
    virtual system by shifting all four colors,
    $(U_gx)(h_1,h_2,h_3,h_4)=x(h_1g,h_2g,h_3g,h_4g)$, the right-regular
    representation on each leg, which the relabeling $h\mapsto h^{-1}$
    turns into the left-regular one
    \cite{gap:rmp_peps_quantum_double_g_isometry}.
\end{definition}

\begin{lemma}[The dual tensor is $G$-isometric]%
    \label{thm:peps_quantum_double_dual_g_isometric}
    \lean{TNLean.PEPS.quantumDoubleDualSpins_mul_colorDiag,
        TNLean.PEPS.quantumDoubleDualSpins_eq_iff,
        TNLean.PEPS.siteMap_quantumDoubleDualTensor_apply,
        TNLean.PEPS.siteMap_quantumDoubleDualTensor_apply_spins,
        TNLean.PEPS.siteMap_quantumDoubleDualTensor_apply_spins_of_mem_invariants,
        TNLean.PEPS.isGInjective_quantumDoubleDualTensor,
        TNLean.PEPS.isGIsometric_quantumDoubleDualTensor,
        TNLean.PEPS.ToricCodeGroup,
        TNLean.PEPS.isGIsometric_toricCodeDualTensor}
    \leanok
    \uses{def:peps_g_injective, def:peps_quantum_double_dual}
    For every finite group $G$ the dual quantum-double tensor is
    $G$-injective and $G$-isometric, with constant $c=|G|$, for the
    color-shift representation
    \cite[Appendix~A, ``The Toric Code and quantum double
        models'']{Cirac2021Matrix},
    \cite[Section~``Examples'', ``The double models'']{Schuch2010PEPS}.
    In particular the toric code tensor, $G=\mathbb{Z}_2$, is
    $\mathbb{Z}_2$-isometric.
\end{lemma}

\begin{proof}\leanok
    Two colorings have the same differences exactly when they differ by one
    global shift $h_i\mapsto h_ig$: the shift of the left color determines
    $g$, and the other three differences then fix the remaining colors.
    Hence $(\mathcal{P}x)(\text{differences of }h)=\sum_{g}x(hg)$. Invariance
    follows by reindexing the sum over colorings by a shift. For an
    invariant $x$ the sum equals $|G|\,x(h)$, so $\mathcal{P}x=0$ forces
    $x=0$, and
    $\langle\mathcal{P}x,\mathcal{P}y\rangle
        =\sum_h\overline{(\mathcal{P}x)(\text{differences of }h)}\,y(h)
        =|G|\langle x,y\rangle$.
\end{proof}

\begin{corollary}[Strings of $R_g$'s on the dual quantum-double network]%
    \label{cor:peps_quantum_double_dual_string_deformation}
    \lean{TNLean.PEPS.rightRegularMatrix,
        TNLean.PEPS.colorShiftRep_eq_torusLegRep,
        TNLean.PEPS.torusBondNetwork_quantumDoubleDual_westSouthOperatorString_eq}
    \leanok
    \uses{thm:peps_quantum_double_dual_g_isometric,
        thm:peps_torus_operator_string_deformation}
    Let $R_g$ be the permutation matrix of $h\mapsto hg$. The color shift of
    the dual tensor is the action of $R_g$ on the incoming and $R_g^{-1}$ on
    the outgoing legs. Hence, in the torus network of the dual tensor with
    colors in $G$ on the bonds, a string of $R_g$'s along the left and bottom
    sides of a rectangle of $m<w$ columns and $n<h$ rows of sites gives the
    same network as the string along its top and right sides
    \cite[Section~``Two dimensions: $G$-injective PEPS'',
        eq.~(2d:move-strings)]{Schuch2010PEPS}.
\end{corollary}

\begin{proof}\leanok
    On the incoming legs $R_g$ sends $x$ to $h\mapsto x(hg)$, and on the
    outgoing legs $R_g^{-1}$ acting along the edge does the same, so the
    action on the four legs is the color shift. The dual tensor is invariant
    under it by Lemma~\ref{thm:peps_quantum_double_dual_g_isometric}, and
    Lemma~\ref{thm:peps_torus_operator_string_deformation} applies.
\end{proof}

Placed at every site of the torus, the dual tensor gives a network whose
bonds carry colors in $G$ and whose physical spins are the differences of the
colors of adjacent bonds; its state is the superposition over all color
patterns \cite[Section~``Examples'', ``The double models'']{Schuch2010PEPS}.
The review identifies it with the equal-weight superposition of the
configurations obeying Gauss' law \cite[Appendix~A, ``The Toric Code and
    quantum double models'']{Cirac2021Matrix}. On the torus not every such
configuration is a difference of colorings: the holonomies of the spins
around the two non-contractible cycles must be trivial.

\begin{lemma}[Parallel sections of edge multipliers on the torus]%
    \label{lem:peps_torus_parallel_section}
    \lean{TNLean.PEPS.zmodTransport, TNLean.PEPS.zmodTransport_zero,
        TNLean.PEPS.zmodTransport_succ,
        TNLean.PEPS.apply_natCast_eq_zmodTransport_mul,
        TNLean.PEPS.zmodTransport_self_eq_one_of_apply_add_one,
        TNLean.PEPS.zmodTransport_val_add_one,
        TNLean.PEPS.IsTorusParallelSection, TNLean.PEPS.IsTorusFlat,
        TNLean.PEPS.IsTorusParallelSection.mul_right,
        TNLean.PEPS.IsTorusParallelSection.isTorusFlat,
        TNLean.PEPS.IsTorusParallelSection.zmodTransport_row,
        TNLean.PEPS.IsTorusParallelSection.zmodTransport_column,
        TNLean.PEPS.IsTorusFlat.zmodTransport_column_mul,
        TNLean.PEPS.IsTorusParallelSection.eq_mul_right,
        TNLean.PEPS.IsTorusFlat.zmodTransport_column_eq_one,
        TNLean.PEPS.exists_isTorusParallelSection_iff,
        TNLean.PEPS.card_isTorusParallelSection}
    \leanok
    \uses{lem:peps_torus_shift_const}
    Let $G$ be a finite group and assign to every site $v$ of the torus two
    elements $a_v,b_v\in G$. Call $f$ from the sites to $G$ a parallel section
    if $f_{v+(1,0)}=a_vf_v$ and $f_{v+(0,1)}=b_vf_v$ for every $v$. Parallel
    sections exist if and only if
    $b_{v+(1,0)}a_v=a_{v+(0,1)}b_v$ for every $v$ and the ordered products
    $a_{(w-1,0)}\cdots a_{(0,0)}$ and $b_{(0,h-1)}\cdots b_{(0,0)}$ are the
    identity. If they exist, there are exactly $|G|$ of them, and any two
    differ by right multiplication with one group element.
\end{lemma}

\begin{proof}\leanok
    If $f$ and $f'$ are parallel sections, $f_v^{-1}f'_v$ is invariant under
    both unit shifts and hence constant, so $f'=fg$; conversely $fg$ is a
    parallel section for every $g$. The two paths from $v$ to $v+(1,1)$ give
    the flatness condition, and going once around a row or a column gives
    the trivial products. Conversely, for flat $a,b$ with trivial products,
    setting $f_{(x,y)}=b_{(x,y-1)}\cdots b_{(x,0)}\,a_{(x-1,0)}\cdots a_{(0,0)}$
    for the representatives $0\le x<w$, $0\le y<h$, flatness moves the
    product up a column one column to the right, conjugated by the
    horizontal multipliers at its ends, so the product around every column is
    trivial; this and the trivial product around the bottom row give the
    two parallel-section rules also where the representatives wrap around.
\end{proof}

\begin{definition}[Regular half-edge coordinates and closed coefficients]%
    \label{def:peps_regular_half_edge_closed}
    \lean{TNLean.PEPS.RegularHalfEdgeConfig,
        TNLean.PEPS.regularHalfEdgeLeftMul,
        TNLean.PEPS.regularHalfEdgeRightMul,
        TNLean.PEPS.regularHalfEdgeOperators,
        TNLean.PEPS.regularHalfEdgeOfBonds,
        TNLean.PEPS.regularHalfEdgeOfOperators,
        TNLean.PEPS.regularHalfEdgeConfigEquivOperatorsBonds,
        TNLean.PEPS.fullRegionHalfEdgeConfigEquiv,
        TNLean.PEPS.regularProjectorClosedState,
        TNLean.PEPS.regularHalfEdgeOfBonds_tail,
        TNLean.PEPS.regularHalfEdgeOfBonds_head,
        TNLean.PEPS.regularHalfEdgeOperators_ofBonds,
        TNLean.PEPS.regularHalfEdgeOperators_ofOperators,
        TNLean.PEPS.regularHalfEdgeOperators_rightMul,
        TNLean.PEPS.regularHalfEdgeOperators_leftMul,
        TNLean.PEPS.regularHalfEdgeOfBonds_eq_rightMul}
    \leanok
    \uses{def:peps_regular_projector_twisted_region}
    Let $G$ be a finite group and let $\Gamma=(V,E)$ be a finite
    simple graph, with each edge oriented from $s(e)$ to $t(e)$.
    A half-edge configuration assigns $\alpha_{v,e}\in G$ to
    every incident vertex-edge pair. Vertex and bond translations are
    $(k\alpha)_{v,e}=k_v\alpha_{v,e}$ and
    $(\alpha r)_{v,e}=\alpha_{v,e}r_e$, respectively.
    Define the operator quotient
    $u(\alpha)_e=\alpha_{t(e),e}\alpha_{s(e),e}^{-1}$.
    For $u,\eta\in G^E$, let $\alpha(u,\eta)$ have tail labels
    $\eta_e$ and head labels $u_e\eta_e$. The map
    $(u,\eta)\mapsto\alpha(u,\eta)$ is a bijection, with inverse
    $\alpha\mapsto(u(\alpha),(\alpha_{s(e),e})_e)$.
    Write $\alpha(u)=\alpha(u,1)$. Then
    $\alpha(u,\eta)=\alpha(u)\eta$,
    $u(\alpha r)=u(\alpha)$, and
    $u(k\alpha)_e=k_{t(e)}u(\alpha)_ek_{s(e)}^{-1}$.

    With $P_v$ the regular averaging projector, the canonical closed
    coefficient vector with inserted operators $u$ is
    \begin{align}
        C_u(\alpha)
        =\sum_{\eta\in G^E}\prod_{v\in V}
        P_v\bigl(\alpha_v,\alpha(u,\eta)_v\bigr).
        \label{eq:peps_half_edge_closed}
    \end{align}
    These are the full-region contractions of
    Definition~\ref{def:peps_regular_projector_twisted_region}.
\end{definition}

\begin{lemma}[Supported expansion of invariant half-edge functions]%
    \label{thm:peps_regular_invariant_state_span}
    \lean{TNLean.PEPS.eq_sum_regularProjectorClosedState_of_invariant,
        TNLean.PEPS.eq_sum_supported_regularProjectorClosedState_of_invariant,
        TNLean.PEPS.mem_span_supported_regularProjectorClosedState_of_invariant,
        TNLean.PEPS.regularProjectorClosedState_eq_zero_of_not_support}
    \leanok
    \uses{def:peps_regular_half_edge_closed}
    Suppose $f(k\alpha)=f(\alpha r)=f(\alpha)$ for every vertex
    translation $k$, bond translation $r$, and half-edge configuration
    $\alpha$. Let $\mathcal P\subseteq G^E$ satisfy
    $u(\alpha)\notin\mathcal P\Rightarrow f(\alpha)=0$. Then
    \begin{align}
        f=\sum_{u\in\mathcal P} f\bigl(\alpha(u)\bigr)C_u.
        \label{eq:peps_supported_orbit_expansion}
    \end{align}
    If $\mathcal P$ is preserved by the endpoint gauges
    $u_e\mapsto k_{t(e)}^{-1}u_ek_{s(e)}$, then every $C_u$ with
    $u\in\mathcal P$ also vanishes when $u(\alpha)\notin\mathcal P$.
    This is the finite regular-coordinate calculation associated with
    the closure argument of \cite[Theorem~5.5]{Schuch2010PEPS}.
\end{lemma}

\begin{proof}\leanok
    \uses{def:peps_regular_half_edge_closed}
    Put $\kappa=|G|$. The independent vertex projectors have kernel
    \begin{align}
        \prod_vP_v(\alpha_v,\beta_v)
        =\kappa^{-|V|}\sum_{k\in G^V}
        \mathbf1_{\{\alpha=k\beta\}}.
        \label{eq:peps_half_edge_average}
    \end{align}
    Summing $f(\beta)$ against (\ref{eq:peps_half_edge_average}) gives
    $\kappa^{-|V|}\sum_k f(k^{-1}\alpha)=f(\alpha)$.
    Reindex $\beta$ by the bijection $(u,\eta)\mapsto\alpha(u,\eta)$.
    Since $\alpha(u,\eta)=\alpha(u)\eta$, right invariance replaces
    $f(\alpha(u,\eta))$ by $f(\alpha(u))$. Summing over $\eta$ gives
    (\ref{eq:peps_half_edge_closed}). If $u\notin\mathcal P$,
    the coefficient $f(\alpha(u))$ is zero because
    $u(\alpha(u))=u$. Removing these zero coefficients gives
    (\ref{eq:peps_supported_orbit_expansion}).

    A nonzero indicator in (\ref{eq:peps_half_edge_average}), with
    $\beta=\alpha(u,\eta)$, implies
    $u(\alpha)_e=k_{t(e)}u_ek_{s(e)}^{-1}$.
    Gauge invariance of $\mathcal P$ therefore forces every indicator
    to vanish outside its prescribed support.
\end{proof}

\begin{lemma}[Physical contraction of the regular closed coefficients]%
    \label{thm:peps_regular_closed_physical_contraction}
    \lean{TNLean.PEPS.regularProjectorClosedState_eq_twistedRegionMatrix,
        TNLean.PEPS.sum_regularProjectorClosedState_eq_stateCoeff}
    \leanok
    \uses{def:peps_regular_half_edge_closed,
        thm:peps_regular_twisted_region_accessibility}
    Let the site maps $A_v$ be regular $G$-injective. For each physical
    configuration $\sigma$, the inserted physical PEPS coefficient is
    \begin{align}
        \sum_\alpha\left(\prod_v A_{v,\sigma_v}(\alpha_v)\right)C_u(\alpha)
        =\sum_{\eta\in G^E}\prod_v
        A_{v,\sigma_v}\bigl(\alpha(u,\eta)_v\bigr).
    \end{align}
\end{lemma}

\begin{proof}\leanok
    \uses{thm:peps_regular_open_region_accessibility}
    Expand (\ref{eq:peps_half_edge_closed}), exchange the two finite sums,
    and factor the sum over independent half-edge labels into one sum
    per vertex. At each vertex use $A_vP_v=A_v$. The remaining sum
    over the shared labels $\eta$ is the inserted contraction.
\end{proof}


\begin{lemma}[Flat torus connections and commuting seams]%
    \label{thm:peps_torus_flat_connection_gauge}
    \lean{TNLean.PEPS.zmodTransport_step,
        TNLean.PEPS.IsTorusFlat.zmodTransport_rectangle,
        TNLean.PEPS.IsTorusFlat.commute_holonomies,
        TNLean.PEPS.torusTreeGauge,
        TNLean.PEPS.IsTorusFlat.torusTreeGauge_right,
        TNLean.PEPS.IsTorusFlat.torusTreeGauge_up,
        TNLean.PEPS.torusNativeRightTransport,
        TNLean.PEPS.torusNativeUpTransport,
        TNLean.PEPS.torusNativeRightTransport_closure,
        TNLean.PEPS.torusNativeUpTransport_closure,
        TNLean.PEPS.torusEdgeAssignment_ext,
        TNLean.PEPS.exists_regularVertexGauge_eq_torusClosure}
    \leanok
    \uses{lem:peps_torus_parallel_section,
        thm:peps_regular_torus_closure_cut_identification}
    Let $G$ be a group and let $a_{x,y},b_{x,y}\in G$ be rightward
    and upward edge elements on $\mathbb Z_w\times\mathbb Z_h$, where
    $w,h>0$. Suppose the connection is flat:
    $b_{x+1,y}a_{x,y}=a_{x,y+1}b_{x,y}$.
    Write $A_m=a_{m-1,0}\cdots a_{0,0}$ and
    $B_{x,n}=b_{x,n-1}\cdots b_{x,0}$, with empty products equal to
    the identity. The holonomies $A_w$ and $B_{0,h}$ commute.
    For the least nonnegative representatives of $x,y$, set
    $k_{x,y}=B_{x,y}A_x$. Then
    \begin{align}
        k_{x+1,y}^{-1}a_{x,y}k_{x,y}
         & =\begin{cases}A_w&x+1=0,\\1&x+1\ne0,\end{cases}
        \label{eq:peps_flat_horizontal_seam}                   \\
        k_{x,y+1}^{-1}b_{x,y}k_{x,y}
         & =\begin{cases}B_{0,h}&y+1=0,\\1&y+1\ne0.\end{cases}
        \label{eq:peps_flat_vertical_seam}
    \end{align}
    If $w,h\ge3$ and the edge elements come from operators $u_e$ on
    the ordered edges of the torus graph, the same gauge transforms
    $u_e$ into the closure assignment with $g=B_{0,h}^{-1}$ and
    horizontal closure element $A_w$. This group-valued seam deformation
    of \cite[Section~5]{Schuch2010PEPS} makes no trivial-holonomy
    assumption.
\end{lemma}

\begin{proof}\leanok
    Induction across a rectangle gives
    \begin{align}
        B_{m,n}A_m
         & =(a_{m-1,n}\cdots a_{0,n})B_{0,n}.
        \label{eq:peps_flat_rectangle_transport}
    \end{align}
    Taking $m=w,n=h$ in
    (\ref{eq:peps_flat_rectangle_transport}) gives
    $B_{0,h}A_w=A_wB_{0,h}$. Before a coordinate wraps, the
    step to its successor adds one factor on the left;
    at the wrap it leaves the full cycle holonomy.
    Applying these two cases to $k$ proves
    (\ref{eq:peps_flat_horizontal_seam}) and
    (\ref{eq:peps_flat_vertical_seam}).
    For the left regular matrices $L_q$, the change on an ordered
    edge $e:s\to t$ is the following contraction, with matrix products
    read from right to left:
    % Formula: L_{k_t^{-1}} L_{u_e} L_{k_s}=L_{k_t^{-1}u_e k_s}.
    % Source: SCP10 paper_v3.tex:1623--1655, oriented string deformation;
    % TorusFlatConnectionGauge.lean, exists_regularVertexGauge_eq_torusClosure.
    % Ink-to-index: L_q(beta,alpha)=1_{beta=q alpha}; the two internal
    % group indices are summed, beta is the output, alpha the input.
    % Boundary signature: (V,P)=(2,0) on both sides. No physical site is drawn.

    \begin{tenkzequation}
        \tenkzkernel
        \begin{tenkz}[rows={wire}, cols=3, bonds=none]
            \tn[at=(1,1), name=flatGauge0,
                ports={180:virtual:$\beta$,0:virtual}]{$L_{k_t^{-1}}$}
            \tn[at=(1,2), name=flatGauge1,
                ports={180:virtual,0:virtual}]{$L_{u_e}$}
            \tn[at=(1,3), name=flatGauge2,
                ports={180:virtual,0:virtual:$\alpha$}]{$L_{k_s}$}
            \tnwire{flatGauge0.0}{flatGauge1.180}
            \tnwire{flatGauge1.0}{flatGauge2.180}
        \end{tenkz}
        $=$
        \begin{tenkz}[rows={wire}, cols=1, bonds=none]
            \tn[at=(1,1), name=flatGauged0,
                ports={180:virtual:$\beta$,0:virtual:$\alpha$}]{$L_{k_t^{-1}u_e k_s}$}
        \end{tenkz}
    \end{tenkzequation}
    On an ordered seam edge the native orientation is reversed.
    Thus its ordered horizontal operator is $A_w^{-1}$ and its
    ordered vertical operator is $B_{0,h}^{-1}$, the
    claimed closure assignment.
\end{proof}

\begin{lemma}[Flat insertions give commuting closure states]%
    \label{thm:peps_flat_insertion_closure_state}
    \lean{TNLean.PEPS.exists_sameState_torusClosure_of_isTorusFlat,
        TNLean.PEPS.exists_stateCoeff_eq_torusGClosure_of_isTorusFlat}
    \leanok
    \uses{thm:peps_torus_flat_connection_gauge,
        thm:peps_regular_closed_gauge,
        thm:peps_regular_torus_closure_cut_identification}
    Let $G$ be finite and let the torus periods be at least three.
    Suppose the site coefficients are invariant under simultaneous
    left translation of their incident group labels. Every flat
    inserted-bond network has the same physical state as a network
    with two commuting closure seams. For a single four-leg tensor,
    this is the closure vector $\psi_{g,h}$ of
    \cite[Section~5]{Schuch2010PEPS}. Neither injectivity nor
    isometry is needed for this equality.
\end{lemma}

\begin{proof}\leanok
    By Lemma~\ref{thm:peps_torus_flat_connection_gauge}, the flat
    connection is gauge-equivalent to the closure assignment, and vertex
    gauges leave the closed contraction unchanged by
    Lemma~\ref{thm:peps_regular_closed_gauge}.
    Lemma~\ref{thm:peps_regular_torus_closure_cut_identification} then
    identifies the resulting ordered-edge contraction with the
    four-leg closure vector.
\end{proof}


\begin{theorem}[Reconstruction from regular inverse coordinates]%
    \label{thm:peps_regular_parent_reconstruction}
    \lean{TNLean.PEPS.globalRegularTensorMap,
        TNLean.PEPS.regularSiteRetraction_eq_self_of_mem_range,
        TNLean.PEPS.globalRegularTensorMap_inverse_eq_self_of_mem_vertexImage,
        TNLean.PEPS.globalRegularTensorMap_inverse_eq_self_of_mem_parent,
        TNLean.PEPS.globalDependentPhysicalMap_regularInverse_leftMul}
    \leanok
    \uses{thm:peps_region_vertex_image_support,
        def:peps_vertex_image_ground_space}
    Let $A_v$ be regular $G$-injective site maps on a finite graph with
    restricted inverses $F_v$ satisfying $F_vA_v=\Pi_v$
    (\cite[Definition~5.1]{Schuch2010PEPS}), where $\Pi_v$ averages
    simultaneous left translation of the incident group labels. Set
    $A_{\mathrm{all}}=\bigotimes_v A_v$ and
    $F_{\mathrm{all}}=\bigotimes_v F_v$.
    If parent regions cover all vertices and $\Psi$ satisfies each
    regional parent condition, then for every $k\in G^V$ (with
    $(k\cdot\alpha)_{v,e}=k_v\alpha_{v,e}$),
    \begin{align}
        A_{\mathrm{all}}F_{\mathrm{all}}\Psi & =\Psi,
        \label{eq:peps_regular_parent_reconstruction}                           \\
        (F_{\mathrm{all}}\Psi)(k\cdot\alpha)
                                             & =(F_{\mathrm{all}}\Psi)(\alpha).
        \label{eq:peps_regular_parent_left_invariance}
    \end{align}
\end{theorem}
\begin{proof}\leanok
    By Theorem~\ref{thm:peps_region_vertex_image_support}, every one-site
    slice of $\Psi$ lies in $\operatorname{im}A_v$, where $A_v\Pi_v=A_v$
    gives $A_vF_vA_v=A_v$:
    % Formula: A_v F_v A_v = A_v Pi_v = A_v, not A_v F_v = I globally.
    % Source: SCP10 Definition 5.1 and paper_v3.tex:1480--1514 (closure inverse).
    % Ink-to-index: matrices read right to left; each virtual wire bundles
    % eta_v in G^{Inc(v)}. A has physical row and virtual column; F the reverse.
    % Joined physical/virtual indices are summed. The external physical row
    % and virtual column remain open: every panel has (V,P)=(1,1).

    \begin{tenkzequation}
        \tenkzkernel
        \begin{tenkz}[rows={wire}, cols=3, bonds=none]
            \tn[at=(1,1), name=parentInverse0,
                ports={180:physical:$i$,0:virtual}]{$A_v$}
            \tn[at=(1,2), name=parentInverse1,
                ports={180:virtual,0:physical}]{$F_v$}
            \tn[at=(1,3), name=parentInverse2,
                ports={180:physical,0:virtual:$\eta_v$}]{$A_v$}
            \tnwire{parentInverse0.0}{parentInverse1.180}
            \tnwire{parentInverse1.0}{parentInverse2.180}
        \end{tenkz}
        $=$
        \begin{tenkz}[rows={wire}, cols=2, bonds=none]
            \tn[at=(1,1), name=parentAverage0,
                ports={180:physical:$i$,0:virtual}]{$A_v$}
            \tn[at=(1,2), name=parentAverage1,
                ports={180:virtual,0:virtual:$\eta_v$}]{$\Pi_v$}
            \tnwire{parentAverage0.0}{parentAverage1.180}
        \end{tenkz}
        $=$
        \begin{tenkz}[rows={wire}, cols=1, bonds=none]
            \tn[at=(1,1), name=parentRetracted0,
                ports={180:physical:$i$,0:virtual:$\eta_v$}]{$A_v$}
        \end{tenkz}
    \end{tenkzequation}
    Tensoring over $v$ gives~(\ref{eq:peps_regular_parent_reconstruction}).
    Applying $F_{\mathrm{all}}$ to~(\ref{eq:peps_regular_parent_reconstruction})
    gives $\phi=F_{\mathrm{all}}\Psi=(\bigotimes_v\Pi_v)\phi$, whose left-translation
    invariance gives~(\ref{eq:peps_regular_parent_left_invariance}).
\end{proof}

\begin{theorem}[Shared-edge invariance and flat support of parent coordinates]%
    \label{thm:peps_regular_parent_coordinate_support}
    \lean{TNLean.PEPS.dependentRegionSlice_regularInverse_halfEdgeRightMul,
        TNLean.PEPS.dependentRegionSlice_regularInverse_apply_eq_zero_of_cycle_ne_one,
        TNLean.PEPS.globalDependentPhysicalMap_regularInverse_singleEdgeRightMul,
        TNLean.PEPS.globalDependentPhysicalMap_regularInverse_rightMul,
        TNLean.PEPS.globalDependentPhysicalMap_regularInverse_eq_zero_of_holonomy_ne_one,
        TNLean.PEPS.globalDependentPhysicalMap_regularInverse_eq_zero_of_torusPlaquetteHolonomy,
        TNLean.PEPS.mem_torusPlaquetteRegion_self,
        TNLean.PEPS.exists_torusPlaquetteRegion_contains_edge,
        TNLean.PEPS.IsGInjective.exists_torusParentCoordinates}
    \leanok
    \uses{thm:peps_regular_parent_reconstruction,
        lem:peps_regular_region_bond_right_invariance,
        thm:peps_regular_region_parent_flatness}
    In the setting of Theorem~\ref{thm:peps_regular_parent_reconstruction},
    if every edge is internal to some parent region, then
    $\phi=F_{\mathrm{all}}\Psi$ satisfies $\phi(\alpha\cdot r)=\phi(\alpha)$
    for all $r\in G^E$, where $(\alpha\cdot r)_{v,e}=\alpha_{v,e}r_e$.
    Orienting each edge from its smaller to its larger endpoint and setting
    $u_e(\alpha)=\alpha_{\mathrm{head}(e),e}\alpha_{\mathrm{tail}(e),e}^{-1}$,
    any closed walk $p$ in a connected parent region with reverse-ordered
    holonomy $\operatorname{Hol}_{u(\alpha)}(p)\ne1$ forces $\phi(\alpha)=0$.
    In particular, on a regular $G$-injective torus with periods at least
    three, the native four-site plaquettes cover all vertices and edges, so
    local inverses $F_v$ exist for which every physical parent ground vector
    satisfies (\ref{eq:peps_regular_parent_reconstruction}),
    (\ref{eq:peps_regular_parent_left_invariance}), right edge invariance,
    and vanishing at nonidentity plaquette holonomies.
\end{theorem}
\begin{proof}\leanok
    Fixing labels outside a parent region $R$ expresses the slice of $\phi$
    as a linear combination of $\bigotimes_{v\in R}F_v$-images of physical
    slices of $\Psi$. By Lemma~\ref{lem:peps_regular_region_bond_right_invariance}
    and Theorem~\ref{thm:peps_regular_region_parent_flatness}, each image is
    invariant under internal-edge right translation and vanishes at
    nonidentity cycle coordinates; translating edges one by one yields
    $\phi(\alpha\cdot r)=\phi(\alpha)$. When all cycle coordinates on a
    connected region $R$ equal $1$, $u_e=q_{\mathrm{head}(e)}q_{\mathrm{tail}(e)}^{-1}$
    has trivial closed-walk holonomy, so $\operatorname{Hol}_{u(\alpha)}(p)\ne1$
    forces $\phi(\alpha)=0$.
\end{proof}

\begin{theorem}[Global spanning of the regular torus parent kernel]%
    \label{thm:peps_regular_torus_parent_spanning}
    \lean{TNLean.PEPS.globalRegularTensorMap_regularProjectorClosedState,
        TNLean.PEPS.IsGInjective.stateCoeff_mem_commutingClosureSpan_of_flat,
        TNLean.PEPS.IsGInjective.torusParentKernel_le_commutingClosureSpan}
    \leanok
    \uses{thm:peps_regular_parent_coordinate_support,
        thm:peps_regular_invariant_state_span,
        thm:peps_regular_closed_physical_contraction,
        thm:peps_flat_insertion_closure_state}
    Let $A$ be regular $G$-injective on a native torus with both periods at
    least three. If $h_v\ge0$ has kernel equal to the open-region ground
    space on the native plaquette $R_v$ and $H=\sum_v h_v\otimes I_{R_v^c}$,
    then (\cite[Theorem~5.7]{Schuch2010PEPS})
    \begin{align}
        \ker H\subseteq
        \operatorname{span}\{\lvert\Psi_{g,h}\rangle:gh=hg\}.
        \label{eq:peps_regular_global_spanning}
    \end{align}
\end{theorem}
\begin{proof}\leanok
    Every $\Psi\in\ker H$ lies in $\ker(h_v\otimes I_{R_v^c})$ for all $v$.
    Let $\phi=F_{\mathrm{all}}\Psi$, let $\mathcal F$ be the set of
    plaquette-flat edge assignments $u$, let $\alpha^u$ have tail label $1$
    and head label $u_e$, and let $C_u$ be the closed contraction of $\Pi_v$
    with insertion $u$. Theorems~\ref{thm:peps_regular_parent_coordinate_support},
    \ref{thm:peps_regular_invariant_state_span},
    and~\ref{thm:peps_regular_closed_physical_contraction} give
    \begin{align}
        \Psi=A_{\mathrm{all}}\phi
         & =\sum_{u\in\mathcal F}\phi(\alpha^u)A_{\mathrm{all}}C_u
        =\sum_{u\in\mathcal F}\phi(\alpha^u)\lvert\Psi_u\rangle.
        \label{eq:peps_regular_flat_expansion}
    \end{align}
    For $\Omega_u=\sum_{\eta\in G^E}\lvert\alpha(u,\eta)\rangle$ with tail
    $\eta_e$ and head $u_e\eta_e$, and $\Pi_{\mathrm{all}}=\bigotimes_v\Pi_v$
    (carrying $|G|^{-|V|}$), we have $C_u=\Pi_{\mathrm{all}}\Omega_u$ and
    $A_{\mathrm{all}}\Pi_{\mathrm{all}}=A_{\mathrm{all}}$:
    % Formula: sum_{u flat} phi(alpha^u) A_all Pi_all Omega_u
    %        = sum_{u flat} phi(alpha^u) A_all Omega_u.
    % Source: SCP10 Theorem 5.7; regular-coordinate proof in
    % TorusRegularParentSpanning.lean, via RegularInvariantStateSpan.lean.
    % Ink-to-index: virtual wire bundles ALL half-edge labels, not one bond;
    % physical wire bundles sigma in the full physical configuration space.
    % Pi_all includes |G|^{-|V|} sum_{k in G^V} L_k; Omega has no scalar factor.
    % All virtual labels contract. Each panel has (V,P)=(0,1), the one
    % physical port denoting the full tuple sigma, not a one-site output.

    \begin{tenkzequation}
        \tenkzkernel
        $\displaystyle\sum_{u\in\mathcal F}\phi(\alpha^u)$
        \begin{tenkz}[rows={wire}, cols=3, bonds=none]
            \tn[at=(1,1), name=parentFlat0,
                ports={180:physical:$\sigma$,0:virtual}]{$A_{\mathrm{all}}$}
            \tn[at=(1,2), name=parentFlat1,
                ports={180:virtual,0:virtual}]{$\Pi_{\mathrm{all}}$}
            \tn[at=(1,3), name=parentFlat2,
                ports={180:virtual}]{$\Omega_u$}
            \tnwire{parentFlat0.0}{parentFlat1.180}
            \tnwire{parentFlat1.0}{parentFlat2.180}
        \end{tenkz}
        $=$
        $\displaystyle\sum_{u\in\mathcal F}\phi(\alpha^u)$
        \begin{tenkz}[rows={wire}, cols=2, bonds=none]
            \tn[at=(1,1), name=parentAbsorbed0,
                ports={180:physical:$\sigma$,0:virtual}]{$A_{\mathrm{all}}$}
            \tn[at=(1,2), name=parentAbsorbed1,
                ports={180:virtual}]{$\Omega_u$}
            \tnwire{parentAbsorbed0.0}{parentAbsorbed1.180}
        \end{tenkz}
    \end{tenkzequation}
    By Theorem~\ref{thm:peps_flat_insertion_closure_state}, each
    $\lvert\Psi_u\rangle$ with $u\in\mathcal F$ equals some commuting
    closure state $\lvert\Psi_{g,h}\rangle$, so
    (\ref{eq:peps_regular_flat_expansion}) yields
    (\ref{eq:peps_regular_global_spanning}).
\end{proof}


\begin{lemma}[Commuting closure labels and centralizers]%
    \label{thm:peps_pair_centralizer_count}
    \lean{TNLean.PEPS.card_commutingPairConjugacyClass_eq_sum_card_centralizer_of_representatives,
        TNLean.PEPS.card_commutingPairConjugacyClass_eq_sum_card_centralizer,
        TNLean.PEPS.commutingPairConjugacyClassEquivSigmaCentralizerOfRepresentatives,
        TNLean.PEPS.commutingPairConjugacyClassEquivSigmaCentralizer,
        TNLean.PEPS.centralizerConjClassEquivCommutingPairFiber}
    \leanok
    \uses{def:peps_pair_conjugacy}
    Let $G$ be a group and choose a representative $g_C$ of each
    conjugacy class $C$. Write $\operatorname{Cl}(H)$ for the set of
    conjugacy classes of a group $H$. The simultaneous-conjugacy classes
    of commuting pairs are in bijection with
    $\coprod_{C\in\operatorname{Cl}(G)}\operatorname{Cl}(Z_G(g_C))$.
    In the group-theoretic classification following
    \cite[Theorem~5.9]{Schuch2010PEPS}, when $G$ is finite their number is
    therefore
    \begin{align}
        K=\sum_{C\in\operatorname{Cl}(G)}
        \bigl|\operatorname{Cl}(Z_G(g_C))\bigr|.
        \label{eq:peps_pair_class_count}
    \end{align}
\end{lemma}

\begin{proof}\leanok
    \uses{thm:peps_pair_conjugacy_centralizer}
    Assign to the class of $(g,h)$ the conjugacy class of $g$.
    If $xgx^{-1}=g_C$, then $xhx^{-1}\in Z_G(g_C)$ and
    $(g,h)$ is simultaneously conjugate to $(g_C,xhx^{-1})$.
    For a fixed first element $g_C$, two such pairs have the same
    simultaneous-conjugacy class exactly when their second elements
    are conjugate in $Z_G(g_C)$, by
    Lemma~\ref{thm:peps_pair_conjugacy_centralizer}.
    This proves the stated bijection. Counting its disjoint components
    gives (\ref{eq:peps_pair_class_count}).
\end{proof}

\begin{lemma}[Dimension of the commuting torus closure span]%
    \label{thm:peps_g_injective_closure_dimension}
    \lean{TNLean.PEPS.IsGInjective.finrank_span_torusGClosure_commuting_eq_sum_card_centralizer,
        TNLean.PEPS.IsGInjective.finrank_span_torusGClosureClass_commuting_eq_sum_card_centralizer,
        TNLean.PEPS.IsGInjective.finrank_span_torusGClosureClass_commuting_of_isSemiRegular,
        TNLean.PEPS.range_torusGClosure_commuting_eq_range_torusGClosureClass}
    \leanok
    \uses{def:peps_semiregular, def:peps_g_injective,
        def:peps_torus_g_closure, def:peps_pair_conjugacy}
    Let $U$ be a semi-regular representation of a finite group $G$,
    and let $A$ be $G$-injective for its four-leg virtual representation.
    On a nonempty rectangular torus, let $\mathcal V_A$ be the span
    of the vectors $\ket{\mathcal M(A\mid C)}$ indexed by commuting
    simultaneous-conjugacy classes $C$. Then the commuting closure span of
    \cite[Theorem~5.9]{Schuch2010PEPS} has dimension
    \begin{align}
        \dim\mathcal V_A
        =\sum_{C\in\operatorname{Cl}(G)}
        \bigl|\operatorname{Cl}(Z_G(g_C))\bigr|.
    \end{align}
\end{lemma}

\begin{proof}\leanok
    \uses{thm:peps_semiregular_g_injective_closure_independence,
        thm:peps_pair_centralizer_count}
    The class-labelled vectors are linearly independent by
    Lemma~\ref{thm:peps_semiregular_g_injective_closure_independence}.
    They form a basis of their span, whose dimension is therefore
    the number of commuting pair classes given by
    Lemma~\ref{thm:peps_pair_centralizer_count}.
\end{proof}


\begin{theorem}[Topology of native plaquettes]%
    \label{thm:peps_torus_plaquette_realization}
    \lean{TNLean.PEPS.torusVertexPoint_add,
        TNLean.PEPS.torusRegionRealization_image_add,
        TNLean.PEPS.connected_torusPlaquetteRegion,
        TNLean.PEPS.torusPlaquetteRegion_eq_image_add,
        TNLean.PEPS.torusPlaquetteRegion_zero_eq_rectangle,
        TNLean.PEPS.isSimplyConnected_torusRegionRealization_plaquette}
    \leanok
    \uses{thm:peps_torus_rectangle_realization}
    On a discrete torus of periods $m,n\ge3$, let
    $R_v=v+\{(0,0),(1,0),(1,1),(0,1)\}$.
    The graph induced on $R_v$ is connected, and the union of its
    closed unit cells in $\mathbb R^2/(m\mathbb Z\times n\mathbb Z)$
    is simply connected, including when $R_v$ crosses a periodic seam.
\end{theorem}
\begin{proof}\leanok
    The four-step plaquette walk has support $R_v$, so its induced graph
    is connected. Translation by $v$ carries the closed-cell realization
    of $R_0$ homeomorphically onto that of $R_v$.
    The realization of $R_0$ is the image of
    $[-\tfrac12,\tfrac32]^2$ under the quotient map.
    Since $2<m,n$, this restriction of the quotient map is injective.
    The compact rectangle is therefore homeomorphic to its image;
    its convexity gives simple connectedness.
\end{proof}

\begin{theorem}[Full regular torus parent ground space]%
    \label{thm:peps_regular_torus_parent_ground_space}
    \lean{TNLean.PEPS.IsGInjective.commutingClosureSpan_le_torusParentKernel,
        TNLean.PEPS.IsGInjective.torusParentKernel_eq_commutingClosureSpan,
        TNLean.PEPS.IsGInjective.finrank_torusParentKernel}
    \leanok
    \uses{thm:peps_torus_plaquette_realization,
        thm:peps_regular_closure_parent_constraint,
        thm:peps_regular_torus_parent_spanning,
        thm:peps_g_injective_closure_dimension}
    Let $A$ be regular $G$-injective on a native torus with both
    periods at least three. For each native plaquette $R_v$, let
    $P_v\ge0$ have kernel equal to the image of the regional PEPS map,
    and let $H=\sum_v(P_v\otimes\mathbf1_{R_v^c})$. Then
    \begin{align}
        \ker H     & =\operatorname{span}\{\lvert A;(g,h)\rangle:gh=hg\},
        \label{eq:peps_regular_parent_kernel}                             \\
        \dim\ker H & =\#\bigl(\{(g,h):gh=hg\}/G\bigr).
        \label{eq:peps_regular_parent_dimension}
    \end{align}
    where $G$ acts by simultaneous conjugation, specializing
    \cite[Theorems~5.7 and~5.9]{Schuch2010PEPS}
    to the regular representation and the indicated torus periods.
\end{theorem}
\begin{proof}\leanok
    Each plaquette is connected with simply connected realization,
    so every commuting closure is annihilated by every $P_v$.
    Conversely, reconstruction from regular inverse coordinates and
    flat-connection reduction express every vector in $\ker H$
    as a linear combination of commuting closures.
    Every site carries $A$. Cut the two seams of the illustrative $3\times3$ torus.
    The coordinate $x$ increases rightward and $y$ increases toward
    the top of the square. Write $\ell_y,r_y$ for the left and right boundary indices and
    $b_x,t_x$ for the bottom and top indices. For $x,y\in\{0,1,2\}$,
    define
    \begin{align}
        B_{g,h}(\ell,r,b,t)
         & =\prod_{y=0}^2 L_h(\ell_y,r_y)
        \prod_{x=0}^2 L_{g^{-1}}(b_x,t_x),
        \label{eq:peps_regular_seam_coefficients}
    \end{align}
    where $L_q(\beta,\alpha)=\mathbf1_{\{\beta=q\alpha\}}$.
    The horizontal native seam label is $h$ and the upward native
    seam label is $g^{-1}$. Reversing the seam edges to the ordered-edge
    convention gives $h^{-1}$ and $g$, respectively. The closure is
    % Formula: |A;(g,h)> = sum_{ell,r,b,t} B_{g,h}(ell,r,b,t)
    % times the open 3x3 patch; gh=hg. Every sum runs independently over G^3.
    % Source: SCP10 paper_v3.tex:1515--1525 (close-AAAA-ug-uh), 1623--1655;
    % TorusFlatConnectionGauge.lean, native right/up closure transport lemmas.
    % Ink-to-index: column x increases rightward, row y increases up the plane;
    % the displayed top row is y=2, bottom row y=0. The left/right/top/bottom
    % exposed virtual legs are ell_y/r_y/t_x/b_x. The nine transverse legs
    % are the physical indices. Unmarked internal bonds contract by identity.
    % Patch boundary signature: (V,P)=(12,9). The four explicit G^3 sums
    % contract all twelve virtual indices against B, giving the full state
    % with (V,P)=(0,9). No normalization or unshown seam contraction is used.
    \begin{tenkzequation}
        \tenkzkernel
        $\displaystyle\lvert A;(g,h)\rangle=
            \sum_{\ell,r,b,t\in G^3}B_{g,h}(\ell,r,b,t)$
        % Tenkz lattice dimensions are literal DSL tokens, not printed products.
        \begin{tenkz}[lattice={3x3}, frame=plane, physical=up] % chktex 29
            \tn[at=(1,1), name=parentSeam,
                ports={180:virtual:$\ell_2$,0:virtual,90:virtual:$t_0$,270:virtual}, label pos=sw]{}
            \tn[at=(1,2), name=seam12,
                ports={180:virtual,0:virtual,90:virtual:$t_1$,270:virtual}, label pos=sw]{}
            \tn[at=(1,3), name=seam13,
                ports={180:virtual,0:virtual:$r_2$,90:virtual:$t_2$,270:virtual}, label pos=sw]{}
            \tn[at=(2,1), name=seam21,
                ports={180:virtual:$\ell_1$,0:virtual,90:virtual,270:virtual}, label pos=sw]{}
            \tn[at=(2,2), name=seam22,
                ports={180:virtual,0:virtual,90:virtual,270:virtual}, label pos=sw]{}
            \tn[at=(2,3), name=seam23,
                ports={180:virtual,0:virtual:$r_1$,90:virtual,270:virtual}, label pos=sw]{}
            \tn[at=(3,1), name=seam31,
                ports={180:virtual:$\ell_0$,0:virtual,90:virtual,270:virtual:$b_0$}, label pos=sw]{}
            \tn[at=(3,2), name=seam32,
                ports={180:virtual,0:virtual,90:virtual,270:virtual:$b_1$}, label pos=sw]{}
            \tn[at=(3,3), name=seam33,
                ports={180:virtual,0:virtual:$r_0$,90:virtual,270:virtual:$b_2$}, label pos=sw]{}
        \end{tenkz}
    \end{tenkzequation}
    Closures in the same simultaneous conjugacy class coincide,
    and one closure per commuting class forms a linearly independent
    family. These statements give both the equality and its dimension.
\end{proof}


\begin{lemma}[The dual quantum-double network is the Gauss-law superposition]%
    \label{thm:peps_quantum_double_dual_gauss_law}
    \lean{TNLean.PEPS.quantumDoubleDualSiteTensor,
        TNLean.PEPS.quantumDoubleDualPEPS,
        TNLean.PEPS.quantumDoubleDualRowTransport,
        TNLean.PEPS.quantumDoubleDualColumnTransport,
        TNLean.PEPS.IsQuantumDoubleGaussLaw,
        TNLean.PEPS.IsQuantumDoubleTrivialHolonomy,
        TNLean.PEPS.isTorusFlat_quantumDoubleDualTransport_iff,
        TNLean.PEPS.forall_eq_quantumDoubleDualSpins_iff,
        TNLean.PEPS.stateCoeff_quantumDoubleDualPEPS}
    \leanok
    \uses{def:peps_quantum_double_dual, lem:peps_torus_site_tensor_contraction,
        lem:peps_torus_parallel_section}
    Let $w,h\ge3$ and place the dual quantum-double tensor at every torus site, with plaquette spins $(s_1,s_2,s_3,s_4)$ (top-left, top-right, bottom-right, bottom-left). A spin configuration obeys Gauss's law if
    \begin{align}
        s_2(v)s_1(v)         & =s_3(v)s_4(v),              &
        s_1(v+(1,0))\,s_2(v) & =s_4(v+(1,1))\,s_3(v+(0,1))
        \label{eq:peps_qd_dual_gauss}
    \end{align}
    at every site $v$ \cite[Section~``Examples'', ``The double models'']{Schuch2010PEPS}, and has trivial holonomy if the ordered products of $a_v=s_2(v)s_1(v)$ along row $y=0$ and of $b_v=s_4(v+(0,1))^{-1}s_1(v)$ along column $x=0$ equal $1\in G$. Then the contracted network coefficient equals $|G|$ on configurations satisfying (\ref{eq:peps_qd_dual_gauss}) with trivial holonomy, and $0$ otherwise \cite{gap:rmp_peps_quantum_double_g_isometry}.
\end{lemma}

\begin{proof}\leanok
    By Lemma~\ref{lem:peps_torus_site_tensor_contraction}, the coefficient counts bond colorings $(f_v,u_v)$ (left and top colors at $v$) matching $(s_1,s_2,s_3,s_4)$:
    \begin{align}
        u_v          & =s_1(v)f_v,    &
        f_{v+(1,0)}  & =a_vf_v,       &
        f_{v+(0,1)}  & =b_vf_v,       &
        s_2(v)s_1(v) & =s_3(v)s_4(v).
        \notag
    \end{align}
    When the first identity of (\ref{eq:peps_qd_dual_gauss}) holds, colorings biject with parallel sections of $(a,b)$, whose flatness is the second identity of (\ref{eq:peps_qd_dual_gauss}); Lemma~\ref{lem:peps_torus_parallel_section} then gives $|G|$ sections for trivial holonomy and $0$ otherwise.
\end{proof}

% Source: SCP10, paper_v3.tex:2625--2923, Section 7.
% Inspected artwork: tc-lattice, tc-blocked-lattice, proj-as-peps,
% tc-tensor-orientation, tc-2x2-renorm and tc-2x2-renorm-syms (figs6 PDFs).
% The unavailable renorm-as-diff-doubles artwork is not reconstructed.
% These are source constructions, not additional formalization claims.
\paragraph{The checkerboard and its blocked plaquettes.}
In the lattice of \cite[Section~7]{Schuch2010PEPS}, each vertex carries
one physical spin in $\C[G]$ and each square is an A or B plaquette. Diagonal
arrows show the source spin orientations; dotted lines record lattice adjacency rather than tensor contractions.
% Source: figs6/tc-lattice.pdf; upper-left A, alternating A/B plaquettes.
% Ink: cb00,...,cb33 are sixteen physical sites, with no tensor ports.
% Diagonal arrows are source spin orientations, not virtual contractions.
% No indices are contracted or opened: this incidence picture has empty
% tensor boundary signature. The spin Hilbert space is C[G] at each dot.
\begin{tenkzequation}
    \tenkzkernel
    \begin{tenkz}[rows={wire,wire,wire,wire,wire,wire,wire},cols=7,bonds=none]
        \tn[at=(1,1),name=cb00,skin=dot]{}
        \tn[at=(1,3),name=cb01,skin=dot]{}
        \tn[at=(1,5),name=cb02,skin=dot]{}
        \tn[at=(1,7),name=cb03,skin=dot]{}
        \tn[at=(3,1),name=cb10,skin=dot]{}
        \tn[at=(3,3),name=cb11,skin=dot]{}
        \tn[at=(3,5),name=cb12,skin=dot]{}
        \tn[at=(3,7),name=cb13,skin=dot]{}
        \tn[at=(5,1),name=cb20,skin=dot]{}
        \tn[at=(5,3),name=cb21,skin=dot]{}
        \tn[at=(5,5),name=cb22,skin=dot]{}
        \tn[at=(5,7),name=cb23,skin=dot]{}
        \tn[at=(7,1),name=cb30,skin=dot]{}
        \tn[at=(7,3),name=cb31,skin=dot]{}
        \tn[at=(7,5),name=cb32,skin=dot]{}
        \tn[at=(7,7),name=cb33,skin=dot]{}
        \tnwire[name=cbH00,stroke=dotted]{cb00}{cb01}
        \tnwire[name=cbH01,stroke=dotted]{cb01}{cb02}
        \tnwire[name=cbH02,stroke=dotted]{cb02}{cb03}
        \tnwire[name=cbH10,stroke=dotted]{cb10}{cb11}
        \tnwire[name=cbH11,stroke=dotted]{cb11}{cb12}
        \tnwire[name=cbH12,stroke=dotted]{cb12}{cb13}
        \tnwire[name=cbH20,stroke=dotted]{cb20}{cb21}
        \tnwire[name=cbH21,stroke=dotted]{cb21}{cb22}
        \tnwire[name=cbH22,stroke=dotted]{cb22}{cb23}
        \tnwire[name=cbH30,stroke=dotted]{cb30}{cb31}
        \tnwire[name=cbH31,stroke=dotted]{cb31}{cb32}
        \tnwire[name=cbH32,stroke=dotted]{cb32}{cb33}
        \tnwire[name=cbV00,stroke=dotted]{cb00}{cb10}
        \tnwire[name=cbV01,stroke=dotted]{cb01}{cb11}
        \tnwire[name=cbV02,stroke=dotted]{cb02}{cb12}
        \tnwire[name=cbV03,stroke=dotted]{cb03}{cb13}
        \tnwire[name=cbV10,stroke=dotted]{cb10}{cb20}
        \tnwire[name=cbV11,stroke=dotted]{cb11}{cb21}
        \tnwire[name=cbV12,stroke=dotted]{cb12}{cb22}
        \tnwire[name=cbV13,stroke=dotted]{cb13}{cb23}
        \tnwire[name=cbV20,stroke=dotted]{cb20}{cb30}
        \tnwire[name=cbV21,stroke=dotted]{cb21}{cb31}
        \tnwire[name=cbV22,stroke=dotted]{cb22}{cb32}
        \tnwire[name=cbV23,stroke=dotted]{cb23}{cb33}
        \tnwire[name=cbArrow00,kind=string,dir=to]{cb00}{0.7 sw of cb00}
        \tnwire[name=cbArrow01,kind=string,dir=to]{cb01}{0.7 nw of cb01}
        \tnwire[name=cbArrow02,kind=string,dir=to]{cb02}{0.7 sw of cb02}
        \tnwire[name=cbArrow03,kind=string,dir=to]{cb03}{0.7 nw of cb03}
        \tnwire[name=cbArrow10,kind=string,dir=to]{cb10}{0.7 se of cb10}
        \tnwire[name=cbArrow11,kind=string,dir=to]{cb11}{0.7 ne of cb11}
        \tnwire[name=cbArrow12,kind=string,dir=to]{cb12}{0.7 se of cb12}
        \tnwire[name=cbArrow13,kind=string,dir=to]{cb13}{0.7 ne of cb13}
        \tnwire[name=cbArrow20,kind=string,dir=to]{cb20}{0.7 sw of cb20}
        \tnwire[name=cbArrow21,kind=string,dir=to]{cb21}{0.7 nw of cb21}
        \tnwire[name=cbArrow22,kind=string,dir=to]{cb22}{0.7 sw of cb22}
        \tnwire[name=cbArrow23,kind=string,dir=to]{cb23}{0.7 nw of cb23}
        \tnwire[name=cbArrow30,kind=string,dir=to]{cb30}{0.7 se of cb30}
        \tnwire[name=cbArrow31,kind=string,dir=to]{cb31}{0.7 ne of cb31}
        \tnwire[name=cbArrow32,kind=string,dir=to]{cb32}{0.7 se of cb32}
        \tnwire[name=cbArrow33,kind=string,dir=to]{cb33}{0.7 ne of cb33}
        \tnmark[name=cbPlaquette00,form=label]{(2,2)}{A}
        \tnmark[name=cbPlaquette01,form=label]{(2,4)}{B}
        \tnmark[name=cbPlaquette02,form=label]{(2,6)}{A}
        \tnmark[name=cbPlaquette10,form=label]{(4,2)}{B}
        \tnmark[name=cbPlaquette11,form=label]{(4,4)}{A}
        \tnmark[name=cbPlaquette12,form=label]{(4,6)}{B}
        \tnmark[name=cbPlaquette20,form=label]{(6,2)}{A}
        \tnmark[name=cbPlaquette21,form=label]{(6,4)}{B}
        \tnmark[name=cbPlaquette22,form=label]{(6,6)}{A}
    \end{tenkz}
\end{tenkzequation}
At a B plaquette, an outward arrow carries $L_u\ket g=\ket{ug}$ and an inward arrow carries $R_u\ket g=\ket{gu^{-1}}$, with four-spin product $V_u$, invariant-space projector $\Pi_B=|G|^{-1}\sum_{u\in G}V_u$, and penalty $h_B=I-\Pi_B$.
For $G=\mathbb Z_2$ ($L_1=R_1=X$, $V_0=I$, $V_1=X^{\otimes4}$),
$\Pi_B=\tfrac12(I+X^{\otimes4})$, $h_B=\tfrac12(I-X^{\otimes4})$, and $h_A=\tfrac12(I-Z^{\otimes4})$ (\cite[Section~7]{Schuch2010PEPS}; the unnormalized source expression $I-\sum_uV_u$ requires the factor $|G|^{-1}$ since $\bigl(\sum_uV_u\bigr)^2=|G|\sum_uV_u$, see \cite{gap:scp10_quantum_double_local_hamiltonian}).
% Source: paper_v3.tex:2640--2656 (binary) and 2870--2872 (general group).
% Normalization correction recorded in
% docs/paper-gaps/scp10_quantum_double_local_hamiltonian.tex.

Grouping the four corner A plaquettes into blocked sites $K_{NW},K_{NE},K_{SE},K_{SW}$ makes the central A plaquette the hole $A_\square$ and the B plaquettes the horizontal and vertical bonds $B_h,B_v$ (shown below after the basis inversion $J$).
% Source: figs6/tc-blocked-lattice.pdf and paper_v3.tex:2808--2844.
% Ink: bk00,...,bk33 are the same sixteen physical sites as cb00,...,cb33.
% Four disjoint enclosures contain the four inner A supports. A_square
% is supported on bk11,bk12,bk22,bk21, one spin from each block.
% B_h uses bk01,bk02,bk12,bk11; B_v uses bk10,bk11,bk21,bk20.
% Spin arrows are reversed from the source by physical basis inversion J.
% The dotted adjacency rails contract nothing; tensor boundary is empty.
\begin{tenkzequation}
    \tenkzkernel
    \begin{tenkz}[rows={wire,wire,wire,wire,wire,wire,wire},cols=7,bonds=none]
        \tn[at=(1,1),name=bk00,skin=dot]{}
        \tn[at=2 e of bk00,name=bk01,skin=dot]{}
        \tn[at=2 e of bk01,name=bk02,skin=dot]{}
        \tn[at=2 e of bk02,name=bk03,skin=dot]{}
        \tn[at=2 s of bk00,name=bk10,skin=dot]{}
        \tn[at=2 e of bk10,name=bk11,skin=dot]{}
        \tn[at=2 e of bk11,name=bk12,skin=dot]{}
        \tn[at=2 e of bk12,name=bk13,skin=dot]{}
        \tn[at=2 s of bk10,name=bk20,skin=dot]{}
        \tn[at=2 e of bk20,name=bk21,skin=dot]{}
        \tn[at=2 e of bk21,name=bk22,skin=dot]{}
        \tn[at=2 e of bk22,name=bk23,skin=dot]{}
        \tn[at=2 s of bk20,name=bk30,skin=dot]{}
        \tn[at=2 e of bk30,name=bk31,skin=dot]{}
        \tn[at=2 e of bk31,name=bk32,skin=dot]{}
        \tn[at=2 e of bk32,name=bk33,skin=dot]{}
        \tnwire[name=bkH00,stroke=dotted]{bk00}{bk01}
        \tnwire[name=bkH01,stroke=dotted]{bk01}{bk02}
        \tnwire[name=bkH02,stroke=dotted]{bk02}{bk03}
        \tnwire[name=bkH10,stroke=dotted]{bk10}{bk11}
        \tnwire[name=bkH11,stroke=dotted]{bk11}{bk12}
        \tnwire[name=bkH12,stroke=dotted]{bk12}{bk13}
        \tnwire[name=bkH20,stroke=dotted]{bk20}{bk21}
        \tnwire[name=bkH21,stroke=dotted]{bk21}{bk22}
        \tnwire[name=bkH22,stroke=dotted]{bk22}{bk23}
        \tnwire[name=bkH30,stroke=dotted]{bk30}{bk31}
        \tnwire[name=bkH31,stroke=dotted]{bk31}{bk32}
        \tnwire[name=bkH32,stroke=dotted]{bk32}{bk33}
        \tnwire[name=bkV00,stroke=dotted]{bk00}{bk10}
        \tnwire[name=bkV01,stroke=dotted]{bk01}{bk11}
        \tnwire[name=bkV02,stroke=dotted]{bk02}{bk12}
        \tnwire[name=bkV03,stroke=dotted]{bk03}{bk13}
        \tnwire[name=bkV10,stroke=dotted]{bk10}{bk20}
        \tnwire[name=bkV11,stroke=dotted]{bk11}{bk21}
        \tnwire[name=bkV12,stroke=dotted]{bk12}{bk22}
        \tnwire[name=bkV13,stroke=dotted]{bk13}{bk23}
        \tnwire[name=bkV20,stroke=dotted]{bk20}{bk30}
        \tnwire[name=bkV21,stroke=dotted]{bk21}{bk31}
        \tnwire[name=bkV22,stroke=dotted]{bk22}{bk32}
        \tnwire[name=bkV23,stroke=dotted]{bk23}{bk33}
        \tnwire[name=bkArrow00,kind=string,dir=to]{bk00}{0.7 ne of bk00}
        \tnwire[name=bkArrow01,kind=string,dir=to]{bk01}{0.7 se of bk01}
        \tnwire[name=bkArrow02,kind=string,dir=to]{bk02}{0.7 ne of bk02}
        \tnwire[name=bkArrow03,kind=string,dir=to]{bk03}{0.7 se of bk03}
        \tnwire[name=bkArrow10,kind=string,dir=to]{bk10}{0.7 nw of bk10}
        \tnwire[name=bkArrow11,kind=string,dir=to]{bk11}{0.7 sw of bk11}
        \tnwire[name=bkArrow12,kind=string,dir=to]{bk12}{0.7 nw of bk12}
        \tnwire[name=bkArrow13,kind=string,dir=to]{bk13}{0.7 sw of bk13}
        \tnwire[name=bkArrow20,kind=string,dir=to]{bk20}{0.7 ne of bk20}
        \tnwire[name=bkArrow21,kind=string,dir=to]{bk21}{0.7 se of bk21}
        \tnwire[name=bkArrow22,kind=string,dir=to]{bk22}{0.7 ne of bk22}
        \tnwire[name=bkArrow23,kind=string,dir=to]{bk23}{0.7 se of bk23}
        \tnwire[name=bkArrow30,kind=string,dir=to]{bk30}{0.7 nw of bk30}
        \tnwire[name=bkArrow31,kind=string,dir=to]{bk31}{0.7 sw of bk31}
        \tnwire[name=bkArrow32,kind=string,dir=to]{bk32}{0.7 nw of bk32}
        \tnwire[name=bkArrow33,kind=string,dir=to]{bk33}{0.7 sw of bk33}
        \tnmark[name=blockNW,form=enclosure,species=selected]{{bk00,bk01,bk10,bk11}}{}
        \tnmark[name=blockNE,form=enclosure,species=selected]{{bk02,bk03,bk12,bk13}}{}
        \tnmark[name=blockSE,form=enclosure,species=selected]{{bk22,bk23,bk32,bk33}}{}
        \tnmark[name=blockSW,form=enclosure,species=selected]{{bk20,bk21,bk30,bk31}}{}
        \tnmark[name=bkPlaquette22,form=label]{(2,2)}{K_{NW}:A}
        \tnmark[name=bkPlaquette26,form=label]{(2,6)}{K_{NE}:A}
        \tnmark[name=bkPlaquette66,form=label]{(6,6)}{K_{SE}:A}
        \tnmark[name=bkPlaquette62,form=label]{(6,2)}{K_{SW}:A}
        \tnmark[name=bkPlaquette44,form=label]{(4,4)}{A_\square}
        \tnmark[name=bkPlaquette24,form=label]{(2,4)}{B_h}
        \tnmark[name=bkPlaquette64,form=label]{(6,4)}{B_h}
        \tnmark[name=bkPlaquette42,form=label]{(4,2)}{B_v}
        \tnmark[name=bkPlaquette46,form=label]{(4,6)}{B_v}
    \end{tenkz}
\end{tenkzequation}

\paragraph{Clockwise physical coordinates.}
With B colors $(p,q,r,s)$ at $(\text{N},\text{E},\text{S},\text{W})$, the original spin values at $(\text{NE},\text{SE},\text{SW},\text{NW})$ are $(qp^{-1},rq^{-1},sr^{-1},ps^{-1})$. Applying the basis involution $J\ket g=\ket{g^{-1}}$ (which satisfies $JL_uJ=R_u$ and $JR_uJ=L_u$, reversing spin arrows) gives the clockwise coordinates
$(a,b,c,d)=(pq^{-1},qr^{-1},rs^{-1},sp^{-1})$.
% Source anchors: paper_v3.tex:2873--2876 (L/R), 2889--2894 (two colors),
% 2896--2908 (clockwise K); original arrows from tc-lattice.pdf.
% This elementary basis comparison is stated mathematically here, without
% identifying it with quantumDoubleDualToK (a different coordinate change).

The A constraint is $abcd=1$ at a K site (clockwise from NE) and $\sigma_{NW,b}\sigma_{SW,a}\sigma_{SE,d}\sigma_{NE,c}=1$ at the hole $A_\square$ (counterclockwise from NW).
In clockwise coordinates on left--right and top--bottom pairs, the horizontal and vertical B-bond actions for $\Pi_B=|G|^{-1}\sum_uV_u$ and $h_B=I-\Pi_B$ are
\begin{align}
    V_u^h & : (a,b,c,d;e,f,g,h)
    \longmapsto(au^{-1},ub,c,d;e,f,gu^{-1},uh), \\
    V_u^v & : (a,b,c,d;e,f,g,h)
    \longmapsto(a,bu^{-1},uc,d;ue,f,g,hu^{-1}).
    \label{eq:peps_source_bond_actions}
\end{align}
% Only the horizontal bond and the northern hole-bond statements below
% carry published declaration links. This source illustration adds none.

\paragraph{A controlled plaquette loop and its two layers.}
For $U_s\in\{L_s,R_s\}$ on a spin adjacent to a B plaquette, the elementary tensor is $P=\sum_{s,k}U_s\ket{k}_{\rm out}\bra{k}_{\rm in}\otimes\ket{s}_v\bra{s}_v$, and contracting four copies around one plaquette gives
\begin{align}
    \operatorname{Loop}(P) & =\sum_{s\in G}
    U_s^{NW}\otimes U_s^{NE}\otimes U_s^{SE}\otimes U_s^{SW}
    =|G|\Pi_B.
    \label{eq:peps_source_p_loop}
\end{align}
% Source: figs6/proj-as-peps.pdf; paper_v3.tex:2696--2722,2879--2885.
% Ink: loopNW,loopNE,loopSE,loopSW are controlled P tensors. Their directed
% virtual edges close clockwise, imposing one common summed control color.
% All physical input/output indices remain open: boundary (virtual,physical)=(0,8).
\begin{tenkzequation}
    \tenkzkernel
    \begin{tenkz}[rows={wire,wire},cols=2,bonds=none]
        \tn[at=(1,1),name=loopNW,skin=box,
            ports={0:virtual,270:virtual,90:physical:$o_1$,180:physical:$i_1$}]{P}
        \tn[at=4 e of loopNW,name=loopNE,skin=box,
            ports={180:virtual,270:virtual,90:physical:$o_2$,0:physical:$i_2$}]{P}
        \tn[at=3 s of loopNE,name=loopSE,skin=box,
            ports={90:virtual,180:virtual,0:physical:$o_3$,270:physical:$i_3$}]{P}
        \tn[at=3 s of loopNW,name=loopSW,skin=box,
            ports={0:virtual,90:virtual,180:physical:$o_4$,270:physical:$i_4$}]{P}
        \tnwire[name=loopN,dir=to]{loopNW.0}{loopNE.180}
        \tnwire[name=loopE,dir=to]{loopNE.270}{loopSE.90}
        \tnwire[name=loopS,dir=to]{loopSE.180}{loopSW.0}
        \tnwire[name=loopW,dir=to]{loopSW.90}{loopNW.270}
        \tnmark[name=loopColor,form=label]{on loopN 0.5}{$s$}
    \end{tenkz}
\end{tenkzequation}

Each spin belongs to two B plaquettes from disjoint layers, whose actions $L_r$ and $R_s$ commute ($L_rR_s=R_sL_r$) and map $\ket1$ to $\ket{rs^{-1}}$ ($\ket{r+s}$ for $G=\mathbb Z_2$). At the southeast spin of a blocked A plaquette in clockwise coordinates, in the original arrow convention this same spin is $sr^{-1}$; applying $J$ gives $rs^{-1}$, defining the four-virtual-leg site tensor $T$:
% Source: paper_v3.tex:2724--2749,2886--2894; tc-tensor-orientation.pdf.
% Ink: lower P^R_s and upper P^L_r have opposite virtual directions.
% One physical wire joins them, and a second joins P^R_s to |1>.
% Result: at original site cb11 / clockwise site bk11, T has virtual
% N,E=r and S,W=s, one physical output rs^{-1} after basis inversion J.
% Both panels have boundary (virtual,physical)=(4,1).
\begin{tenkzequation}
    \tenkzkernel
    \begin{tenkz}[rows={wire,wire,wire},cols=1,bonds=none]
        \tn[at=(1,1),name=stackL,skin=box,
            ports={135:virtual,45:virtual,90:physical:$rs^{-1}$,270:physical}]{P^L}
        \tn[at=2 s of stackL,name=stackR,skin=box,
            ports={225:virtual,315:virtual,90:physical,270:physical}]{P^R}
        \tn[at=2 s of stackR,name=stackOne,skin=box,ports={90:physical}]{\ket1}
        \tnwire[name=stackLIn,dir=from]{stackL.135}{open nw}
        \tnwire[name=stackLOut,dir=to]{stackL.45}{open ne}
        \tnwire[name=stackRIn,dir=from]{stackR.225}{open sw}
        \tnwire[name=stackROut,dir=to]{stackR.315}{open se}
        \tnmark[form=label,label pos=w]{on stackLIn 0.5}{$r$}
        \tnmark[form=label,label pos=e]{on stackLOut 0.5}{$r$}
        \tnmark[form=label,label pos=w]{on stackRIn 0.5}{$s$}
        \tnmark[form=label,label pos=e]{on stackROut 0.5}{$s$}
        \tnwire[name=stackMiddle,dir=to]{stackR.90}{stackL.270}
        \tnwire[name=stackInput,dir=to]{stackOne.90}{stackR.270}
    \end{tenkz}
    \quad = \quad
    \begin{tenkz}[rows={wire},cols=1,bonds=none]
        \tn[at=(1,1),name=stackT,skin=box,
            ports={90:virtual,0:virtual,270:virtual,180:virtual,
                    135:physical:$rs^{-1}$}]{T}
        \tnwire[name=stackTNorth,dir=from]{stackT.90}{open n}
        \tnwire[name=stackTEast,dir=to]{stackT.0}{open e}
        \tnwire[name=stackTSouth,dir=to]{stackT.270}{open s}
        \tnwire[name=stackTWest,dir=from]{stackT.180}{open w}
        \tnmark[form=label,label pos=e]{on stackTNorth 0.5}{$r$}
        \tnmark[form=label,label pos=n]{on stackTEast 0.5}{$r$}
        \tnmark[form=label,label pos=e]{on stackTSouth 0.5}{$s$}
        \tnmark[form=label,label pos=n]{on stackTWest 0.5}{$s$}
    \end{tenkz}
\end{tenkzequation}

\paragraph{The source renormalization and the open boundary identity.}
Grouping four $T$ tensors around an A plaquette via the pair permutation $\ket{a,b}\mapsto\ket{a,b^{-1}a}$ (or CNOT $\ket{a,b}\mapsto\ket{a,a+b}$ for $G=\mathbb Z_2$, conjugating $X\otimes X\mapsto X\otimes I$ and $Z\otimes Z\mapsto I\otimes Z$) separates $K$ from one-leg factors $L$, as in the binary source picture below. In the explicit copied-color coordinates of Theorem~\ref{thm:peps_qd_physical_renorm}, the four surplus virtual factors are $\bra1$; they are not invariant under regular translations. In a basis where a one-leg factor instead satisfies $LU_u=L$ for a regular representation ($u\in G$), the identity $|G|^{-1}\sum_{u\in G}U_u=\ket\omega\bra\omega$ gives the conditional factorization
\begin{align}
    L                        & =L\,|G|^{-1}\sum_{u\in G}U_u=L\ket\omega\bra\omega, &
    \sum_kL\ket k\otimes M_k & =L\ket\omega\otimes
    \sum_k\braket{\omega}{k}M_k, \notag
\end{align}
so in that regular-invariant basis the physical output factors off as $L\ket\omega$. The diagram depicts the separated factors without specifying their physical dimensions:
% Source: paper_v3.tex:2777--2806 and figs6/tc-2x2-renorm.pdf.
% Ink: sourceK carries four virtual indices and a bundled physical output;
% each sourceL* carries ONE virtual index AND a distinct physical output.
% K virtual arrows enter at N,W and leave at E,S; every L arrow enters,
% as in tc-2x2-renorm.pdf. No factors are connected in this block drawing.
% Boundary (virtual,physical)=(8,5); physical ports denote output factors,
% not five elementary qubits and not the four fixed virtual CNOT registers.
% Source physical RG statement only: no leanok or declaration is attached.
\begin{tenkzequation}
    \tenkzkernel
    \begin{tenkz}[rows={wire,wire,wire},cols=3,bonds=none]
        \tn[at=(2,2),name=sourceK,skin=box,
            ports={90:virtual,0:virtual,270:virtual,180:virtual,135:physical:$\mathcal H_K$}]{K}
        \tn[at=3 n of sourceK,name=sourceLN,skin=box,
            ports={90:virtual,0:physical:$\mathcal H_{L_N}$}]{L_N}
        \tn[at=3 e of sourceK,name=sourceLE,skin=box,
            ports={0:virtual,90:physical:$\mathcal H_{L_E}$}]{L_E}
        \tn[at=3 s of sourceK,name=sourceLS,skin=box,
            ports={270:virtual,0:physical:$\mathcal H_{L_S}$}]{L_S}
        \tn[at=3 w of sourceK,name=sourceLW,skin=box,
            ports={180:virtual,270:physical:$\mathcal H_{L_W}$}]{L_W}
        \tnwire[name=sourceKn,dir=from]{sourceK.90}{open n}
        \tnwire[name=sourceKe,dir=to]{sourceK.0}{open e}
        \tnwire[name=sourceKs,dir=to]{sourceK.270}{open s}
        \tnwire[name=sourceKw,dir=from]{sourceK.180}{open w}
        \tnwire[name=sourceLNn,dir=from]{sourceLN.90}{open n}
        \tnwire[name=sourceLEe,dir=from]{sourceLE.0}{open e}
        \tnwire[name=sourceLSs,dir=from]{sourceLS.270}{open s}
        \tnwire[name=sourceLWw,dir=from]{sourceLW.180}{open w}
    \end{tenkz}
\end{tenkzequation}

At the coefficient level on the original four physical spins and eight virtual boundary coordinates, the binary pair CNOTs satisfy $\mathsf B\mathsf C=\mathsf K E^\dagger$, where $E$ sets the four second virtual registers to zero (Lemma~\ref{thm:peps_kitaev_checkerboard_map_identity}).


% SCP10, paper_v3.tex:2832--2940, especially 2861--2877 and 2896--2923.
% The available tc-lattice and tc-blocked-lattice artwork was inspected.
% The unavailable renorm-as-diff-doubles artwork is not reconstructed here.
% Declaration links are added only after the corresponding statements pass.
\paragraph{The clockwise color-difference tensor.}
Let $G$ be a finite group. At one blocked site, place virtual colors $(p,q,r,s)$ on the $(\text{N},\text{E},\text{S},\text{W})$ legs and physical spins $(a,b,c,d)$ on the $(\text{NE},\text{SE},\text{SW},\text{NW})$ corners, following the clockwise convention of \cite[Equation~(7.10) and footnote]{Schuch2010PEPS}.

\begin{definition}[The clockwise blocked tensor]
    \label{def:peps_quantum_double_k_tensor}
    \lean{TNLean.PEPS.quantumDoubleKSpins,TNLean.PEPS.quantumDoubleKTensor,
        TNLean.PEPS.siteMap_quantumDoubleKTensor_apply}
    \leanok
    Define the color-difference map $F:G^4\to G^4$, the unnormalized tensor $K$, and its linear map $\mathsf K\ket{p,q,r,s}=\ket{F(p,q,r,s)}$ by
    \begin{align}
        F(p,q,r,s)         & =(pq^{-1},qr^{-1},rs^{-1},sp^{-1}), \\
        K(p,q,r,s;a,b,c,d) & =\delta_{(a,b,c,d),F(p,q,r,s)}.
        \label{eq:peps_qd_k_tensor}
    \end{align}
\end{definition}
% Formula: K(p,q,r,s;a,b,c,d)=delta_((a,b,c,d),F(p,q,r,s)).
% Ink: virtual N,E,S,W = p,q,r,s; physical NE,SE,SW,NW = a,b,c,d.
% All eight legs are open; no repeated label is summed in this picture.
% External boundary signature: (virtual,physical)=(4,4), alphabet G per leg.
\begin{tenkzequation}
    \tenkzkernel
    \begin{tenkz}[rows={wire},cols=1,bonds=none]
        \tn[at=(1,1),name=qdK,skin=box,
            ports={90:virtual:$p$,0:virtual:$q$,270:virtual:$r$,180:virtual:$s$,
                    45:physical:$a$,315:physical:$b$,225:physical:$c$,135:physical:$d$}]{K}
    \end{tenkz}
\end{tenkzequation}

\begin{theorem}[Comparison of physical orientation conventions]
    \label{thm:peps_quantum_double_k_orientation}
    \lean{TNLean.PEPS.quantumDoubleDualToK,
        TNLean.PEPS.quantumDoubleDualToK_spins,
        TNLean.PEPS.quantumDoubleKTensor_eq_dual}
    \leanok
    \uses{def:peps_quantum_double_k_tensor}
    For $F_{\rm dual}(p,q,r,s)=(ps^{-1},qp^{-1},qr^{-1},rs^{-1})$ and $J(y_1,y_2,y_3,y_4)=(y_2^{-1},y_3,y_4,y_1^{-1})$, one has $J\circ F_{\rm dual}=F$.
\end{theorem}
\begin{proof}\leanok
    Since $(qp^{-1})^{-1}=pq^{-1}$ and $(ps^{-1})^{-1}=sp^{-1}$, $J(F_{\rm dual}(p,q,r,s))=F(p,q,r,s)$.
\end{proof}
% J is a local physical coordinate permutation of two blocked conventions;
% it is not a unitary renormalization of the unblocked tensor T.

\begin{definition}[The local flatness projector]
    \label{def:peps_quantum_double_local_projector}
    \lean{TNLean.PEPS.quantumDoubleKHolonomy,
        TNLean.PEPS.quantumDoubleKLocalProjector,TNLean.PEPS.quantumDoubleKLocalTerm,
        TNLean.PEPS.quantumDoubleKLocalProjector_mulVec}
    \leanok
    \uses{def:peps_quantum_double_k_tensor}
    On $\C[G^4]$ (for $G$ not necessarily abelian), set
    \begin{align}
        P_0 & =\sum_{abcd=1}\ketbra{a,b,c,d}{a,b,c,d},
            & h_0                                      & =I-P_0.
        \label{eq:peps_qd_local_projector}
    \end{align}
\end{definition}

\begin{theorem}[Local tensor image]
    \label{thm:peps_quantum_double_local_image}
    \lean{TNLean.PEPS.quantumDoubleKSpins_range,
        TNLean.PEPS.quantumDoubleKHolonomy_spins,
        TNLean.PEPS.quantumDoubleKColors,TNLean.PEPS.quantumDoubleKSpins_colors,
        TNLean.PEPS.siteMap_quantumDoubleKTensor_single,
        TNLean.PEPS.quantumDoubleKLocalProjector_isStarProjection,
        TNLean.PEPS.quantumDoubleKLocalTerm_isStarProjection,
        TNLean.PEPS.quantumDoubleKLocalProjector_siteMap,
        TNLean.PEPS.range_siteMap_quantumDoubleKTensor,
        TNLean.PEPS.range_siteMap_quantumDoubleKTensor_eq_ker_localTerm}
    \leanok
    \uses{def:peps_quantum_double_local_projector}
    The operators $P_0$ and $h_0$ are orthogonal projectors satisfying $P_0\mathsf K=\mathsf K$, $\operatorname{ran}\mathsf K=\ker h_0$, and
    \begin{align}
        \operatorname{ran}\mathsf K
         & =\operatorname{ran}P_0
        =\operatorname{span}\{\ket{a,b,c,d}:abcd=1\}.
        \label{eq:peps_qd_local_image}
    \end{align}
\end{theorem}
\begin{proof}\leanok
    Every $(a,b,c,d)=F(p,q,r,s)$ satisfies $abcd=1$, and conversely $abcd=1$ gives $F(1,a^{-1},(ab)^{-1},(abc)^{-1})=(a,b,c,d)$, proving (\ref{eq:peps_qd_local_image}). Since $P_0$ is diagonal with $\{0,1\}$ entries, $P_0$ and $h_0=I-P_0$ are orthogonal projectors and $\operatorname{ran}P_0=\ker h_0$.
\end{proof}

\paragraph{One shared bond.}
Contracting the east leg of a left tensor with the west leg of a right tensor over $x\in G$ at fixed exterior colors $(p,r,s;t,v,w)$ gives
\begin{align}
    \Psi_{p,r,s;t,v,w}
     & =\sum_{x\in G}\ket{F(p,x,r,s)}\otimes\ket{F(t,v,w,x)}.
    \label{eq:peps_qd_bond_contraction}
\end{align}
% Formula: Psi=sum_x |F(p,x,r,s)> tensor |F(t,v,w,x)>.
% Ink: the sole internal virtual wire is q(left)=s(right)=x, summed once.
% The six other virtual colors p,r,s,t,v,w are open boundary arguments.
% Each physical leg is a complete G^4 tuple, not one selected corner spin.
% External boundary signature: (virtual,physical)=(6,2).
\begin{tenkzequation}
    \tenkzkernel
    \begin{tenkz}[rows={wire},cols=2,bonds=none]
        \tn[at=(1,1),name=qdLeft,skin=box,
            ports={90:virtual:$p$,0:virtual,270:virtual:$r$,180:virtual:$s$,
                    135:physical:$\sigma_L$}]{K_L}
        \tn[at=3 e of qdLeft,name=qdRight,skin=box,
            ports={90:virtual:$t$,0:virtual:$v$,270:virtual:$w$,180:virtual,
                    45:physical:$\sigma_R$}]{K_R}
        \tnwire[name=qdShared]{qdLeft.0}{qdRight.180}
        \tnmark[form=label,label pos=n]{on qdShared 0.5}{$x$}
    \end{tenkz}
\end{tenkzequation}

\begin{definition}[Coherent action at an east--west bond]
    \label{def:peps_quantum_double_bond_average}
    \lean{TNLean.PEPS.quantumDoubleKBondContraction,
        TNLean.PEPS.quantumDoubleKBondContraction_eq,
        TNLean.PEPS.quantumDoubleKBondLeft,TNLean.PEPS.quantumDoubleKBondRight,
        TNLean.PEPS.quantumDoubleKBondAction,TNLean.PEPS.quantumDoubleKBondPermutation,
        TNLean.PEPS.quantumDoubleKBondMatrix,TNLean.PEPS.quantumDoubleKBondMatrix_mulVec,
        TNLean.PEPS.quantumDoubleKBondAverage,
        TNLean.PEPS.quantumDoubleKBondTerm}
    \leanok
    \uses{def:peps_quantum_double_k_tensor}
    For $u\in G$, define $V_u,A_e,h_e$ on $\C[G^4]\otimes\C[G^4]$ (\cite[Section~7, ``The double models'']{Schuch2010PEPS}) by
    \begin{align}
        V_u\bigl(\ket{a,b,c,d}\otimes\ket{e,f,g,h}\bigr)
            & =\ket{au^{-1},ub,c,d}\otimes\ket{e,f,gu^{-1},uh},           \\
        A_e & =\frac1{|G|}\sum_{u\in G}V_u,
            & h_e                                               & =I-A_e.
        \label{eq:peps_qd_bond_average}
    \end{align}
\end{definition}

\begin{theorem}[The normalized bond projector fixes the contraction]
    \label{thm:peps_quantum_double_bond_fixed}
    \lean{TNLean.PEPS.quantumDoubleKBondAction_mul,
        TNLean.PEPS.quantumDoubleKBondAction_one,
        TNLean.PEPS.quantumDoubleKSpins_bondLeft,TNLean.PEPS.quantumDoubleKSpins_bondRight,
        TNLean.PEPS.quantumDoubleKBondMatrix_single,
        TNLean.PEPS.quantumDoubleKBondAverage_isStarProjection,
        TNLean.PEPS.quantumDoubleKBondTerm_isStarProjection,
        TNLean.PEPS.quantumDoubleKBondAverage_mulVec_eq_self_iff,
        TNLean.PEPS.quantumDoubleKBondContraction_invariant,
        TNLean.PEPS.quantumDoubleKBondAverage_contraction}
    \leanok
    \uses{def:peps_quantum_double_bond_average}
    The operators satisfy $V_uV_v=V_{uv}$ and $V_u^\dagger=V_{u^{-1}}$, so $A_e$ and $h_e$ are orthogonal projectors, $A_ev=v\iff\forall u\in G,\,V_uv=v$, and for all $(p,r,s;t,v,w)$,
    \begin{align}
        V_u\Psi_{p,r,s;t,v,w} & =\Psi_{p,r,s;t,v,w},        \\
        A_e\Psi_{p,r,s;t,v,w} & =\Psi_{p,r,s;t,v,w},
                              & h_e\Psi_{p,r,s;t,v,w} & =0.
        \label{eq:peps_qd_bond_fixed}
    \end{align}
\end{theorem}
\begin{proof}\leanok
    Since $V_u$ is a unitary representation of $G$, $A_e=|G|^{-1}\sum_{u\in G}V_u$ is the orthogonal projector onto the $G$-invariant subspace. In (\ref{eq:peps_qd_bond_contraction}), $V_u$ sends the summand at $x$ to the summand at $ux$, so reindexing $x\mapsto ux$ gives (\ref{eq:peps_qd_bond_fixed}).
\end{proof}

The factor $|G|^{-1}$ in (\ref{eq:peps_qd_bond_average}) is necessary: for $S_e=\sum_uV_u$, $S_e^2=|G|S_e$ and $(I-S_e)\psi=(1-|G|)\psi$ on invariant vectors, so the unnormalized sum in \cite[Section~7]{Schuch2010PEPS} must be divided by $|G|$.
% Local correction and remaining scope are recorded in
% docs/paper-gaps/scp10_quantum_double_local_hamiltonian.tex.

\begin{theorem}[Two-block tensor image]
    \label{thm:peps_quantum_double_bond_image}
    \lean{TNLean.PEPS.quantumDoubleKPairLocalProjector,
        TNLean.PEPS.quantumDoubleKHolonomy_bondLeft,
        TNLean.PEPS.quantumDoubleKHolonomy_bondRight,
        TNLean.PEPS.quantumDoubleKPairLocalProjector_isStarProjection,
        TNLean.PEPS.quantumDoubleKPairLocalProjector_commute_bondAverage,
        TNLean.PEPS.quantumDoubleKBondSupportProjector,
        TNLean.PEPS.quantumDoubleKBondSupportProjector_isStarProjection,
        TNLean.PEPS.quantumDoubleKPairLocalProjector_contraction,
        TNLean.PEPS.quantumDoubleKBondSupportProjector_contraction,
        TNLean.PEPS.quantumDoubleKBondBoundaryColors,
        TNLean.PEPS.quantumDoubleKBondBoundaryColors_spins,
        TNLean.PEPS.quantumDoubleKBondAverage_single_eq_contraction,
        TNLean.PEPS.range_quantumDoubleKBondContraction,
        TNLean.PEPS.mem_range_quantumDoubleKBondContraction_iff}
    \leanok
    \uses{def:peps_quantum_double_local_projector,
        thm:peps_quantum_double_bond_fixed}
    Let $\mathsf C\ket\beta=\Psi_\beta$ be the contraction map of (\ref{eq:peps_qd_bond_contraction}) and set $P_{00}=P_0\otimes P_0$. Then $P_{00}A_e=A_eP_{00}$, $Q_e=A_eP_{00}$ is an orthogonal projector, and
    \begin{align}
        \operatorname{ran}\mathsf C & =\operatorname{ran}Q_e
        =\{v:P_{00}v=v\text{ and }\forall u\in G,\,V_uv=v\}.
        \label{eq:peps_qd_bond_image}
    \end{align}
\end{theorem}
\begin{proof}\leanok
    Each $V_u$ preserves $abcd$ and $efgh$, so $V_u$ and $A_e$ commute with $P_{00}$, making $Q_e=A_eP_{00}$ an orthogonal projector with $\operatorname{ran}\mathsf C\subseteq\operatorname{ran}Q_e$. Conversely, for $\xi=(a,b,c,d;e,f,g,h)$ with $abcd=efgh=1$, setting $\beta(\xi)=(a,b^{-1},(bc)^{-1};efg,fg,g)$ gives
    \begin{align}
        A_e\ket\xi & =|G|^{-1}\Psi_{\beta(\xi)}\in\operatorname{ran}\mathsf C,
    \end{align}
    while $P_{00}\ket\xi=0$ otherwise, so $\operatorname{ran}Q_e\subseteq\operatorname{ran}\mathsf C$.
\end{proof}

\paragraph{The hole between four blocked tensors.}
Contracting four copies of $K$ at $(\text{NW},\text{NE},\text{SE},\text{SW})$ over the shared $(\text{N},\text{W},\text{S},\text{E})$ bond colors $(x,y,z,w)\in G^4$ at fixed exterior colors $\beta=(p,s,t,v,k,l,m,n)$ gives
\begin{align}
    \Phi_\beta
     & =\sum_{x,y,z,w\in G}
    \ket{F(p,x,y,s)}\otimes\ket{F(t,v,w,x)}\nonumber                 \\
     & \hspace{2.5em}\otimes\ket{F(w,k,l,z)}\otimes\ket{F(y,z,m,n)}.
    \label{eq:peps_qd_hole_contraction}
\end{align}
% Formula: Phi_beta is the literal four-K sum in eq:peps_qd_hole_contraction.
% Ink: NW.E=NE.W=x; NW.S=SW.N=y; SW.E=SE.W=z; NE.S=SE.N=w.
% The four internal virtual colors are each summed once; exterior colors
% p,s,t,v,k,l,m,n are independent, unsummed boundary arguments.
% Physical sigma_NW,sigma_NE,sigma_SE,sigma_SW are complete G^4 tuples.
% The selected hole spins are NW.b,SW.a,SE.d,NE.c, in that product order.
% External boundary signature: (virtual,physical)=(8,4).
\begin{tenkzequation}
    \tenkzkernel
    \begin{tenkz}[rows={wire,wire},cols=2,bonds=none]
        \tn[at=(1,1),name=qdNW,skin=box,
            ports={90:virtual:$p$,0:virtual,270:virtual,180:virtual:$s$,
                    135:physical:$\sigma_{NW}$}]{K}
        \tn[at=3 e of qdNW,name=qdNE,skin=box,
            ports={90:virtual:$t$,0:virtual:$v$,270:virtual,180:virtual,
                    45:physical:$\sigma_{NE}$}]{K}
        \tn[at=3 s of qdNE,name=qdSE,skin=box,
            ports={90:virtual,0:virtual:$k$,270:virtual:$l$,180:virtual,
                    315:physical:$\sigma_{SE}$}]{K}
        \tn[at=3 s of qdNW,name=qdSW,skin=box,
            ports={90:virtual,0:virtual,270:virtual:$m$,180:virtual:$n$,
                    225:physical:$\sigma_{SW}$}]{K}
        \tnwire[name=qdNorth]{qdNW.0}{qdNE.180}
        \tnwire[name=qdWest]{qdNW.270}{qdSW.90}
        \tnwire[name=qdSouth]{qdSW.0}{qdSE.180}
        \tnwire[name=qdEast]{qdNE.270}{qdSE.90}
        \tnmark[form=label,label pos=n]{on qdNorth 0.5}{$x$}
        \tnmark[form=label,label pos=w]{on qdWest 0.5}{$y$}
        \tnmark[form=label,label pos=s]{on qdSouth 0.5}{$z$}
        \tnmark[form=label,label pos=e]{on qdEast 0.5}{$w$}
    \end{tenkz}
\end{tenkzequation}

\begin{definition}[The four-block plaquette projector]
    \label{def:peps_quantum_double_hole_projector}
    \lean{TNLean.PEPS.QuantumDoubleKPlaquetteConfig,
        TNLean.PEPS.quantumDoubleKPlaquetteSpins,
        TNLean.PEPS.quantumDoubleKPlaquetteContraction,
        TNLean.PEPS.quantumDoubleKPlaquetteHolonomy,
        TNLean.PEPS.quantumDoubleKPlaquetteProjector,
        TNLean.PEPS.quantumDoubleKPlaquetteTerm}
    \leanok
    \uses{def:peps_quantum_double_k_tensor}
    For $\sigma=(\sigma_{NW},\sigma_{NE},\sigma_{SE},\sigma_{SW})\in(G^4)^4$, define on $\C[(G^4)^4]$
    \begin{align}
        H_\square(\sigma)   & =\sigma_{NW,b}\sigma_{SW,a}
        \sigma_{SE,d}\sigma_{NE,c},                                                    \\
        P_\square\ket\sigma & =\delta_{H_\square(\sigma),1}\ket\sigma,
                            & h_\square                                & =I-P_\square.
        \label{eq:peps_qd_hole_projector}
    \end{align}
\end{definition}

\begin{theorem}[Plaquette flatness of the four-block contraction]
    \label{thm:peps_quantum_double_hole_fixed}
    \lean{TNLean.PEPS.quantumDoubleKPlaquetteHolonomy_spins,
        TNLean.PEPS.quantumDoubleKPlaquetteProjector_isStarProjection,
        TNLean.PEPS.quantumDoubleKPlaquetteTerm_isStarProjection,
        TNLean.PEPS.quantumDoubleKPlaquetteContraction_eq,
        TNLean.PEPS.quantumDoubleKPlaquetteProjector_contraction}
    \leanok
    \uses{def:peps_quantum_double_hole_projector}
    The operators $P_\square$ and $h_\square$ are orthogonal projectors, and $P_\square\Phi_\beta=\Phi_\beta$, $h_\square\Phi_\beta=0$ for every $\beta$.
\end{theorem}
\begin{proof}\leanok
    Each summand of (\ref{eq:peps_qd_hole_contraction}) has hole spins $(xy^{-1},yz^{-1},zw^{-1},wx^{-1})$ with product $(xy^{-1})(yz^{-1})(zw^{-1})(wx^{-1})=1$, so the $\{0,1\}$-diagonal projector $P_\square$ fixes $\Phi_\beta$.
\end{proof}

\begin{theorem}[Commutation at the northern bond]
    \label{thm:peps_quantum_double_hole_commute}
    \lean{TNLean.PEPS.quantumDoubleKPlaquetteTopBond,
        TNLean.PEPS.quantumDoubleKPlaquetteHolonomy_topBond,
        TNLean.PEPS.quantumDoubleKPlaquetteHolonomy_topBond_eq_one_iff,
        TNLean.PEPS.quantumDoubleKPlaquetteProjector_commute_topBond,
        TNLean.PEPS.quantumDoubleKPlaquetteProjector_commute_topBondAverage}
    \leanok
    \uses{def:peps_quantum_double_bond_average,
        def:peps_quantum_double_hole_projector}
    Let $W_u=V_u\otimes I$ act by $V_u$ on $(\sigma_{NW},\sigma_{NE})$ and the identity on $(\sigma_{SE},\sigma_{SW})$. Then $P_\square$ commutes with $W_u$ for all $u\in G$ and with $|G|^{-1}\sum_{u\in G}W_u$.
\end{theorem}
\begin{proof}\leanok
    Since $H_\square(W_u\sigma)=uH_\square(\sigma)u^{-1}$, one has $H_\square(W_u\sigma)=1\iff H_\square(\sigma)=1$, so $P_\square$ commutes with each $W_u$ and with their average.
\end{proof}

% Scope: exact local and east--west two-block terms, and a four-block hole.
% These statements do not identify a full-lattice common kernel, prove
% commutation of all placed terms, or construct the unblocked T-to-K unitary.


\begin{definition}[Primal quantum-double tensor]%
    \label{def:peps_quantum_double_primal}
    \lean{TNLean.PEPS.quantumDoublePrimalLabels,
        TNLean.PEPS.quantumDoublePrimalTensor}
    \leanok
    The primal tensor $A^{g_1,g_2,g_3,g_4}_{t,r,b,l}$ with plaquette spins $(g_1,g_2,g_3,g_4)$ (top-left, top-right, bottom-right, bottom-left) equals $1$ when
    \begin{align}
        (t,r,b,l) & =(g_1g_2,\,g_2^{-1}g_3,\,g_4g_3,\,g_1g_4^{-1}),
        \label{eq:peps_qd_primal_labels}
    \end{align}
    and $0$ otherwise \cite[Appendix~A, ``The Toric Code and quantum double models'']{Cirac2021Matrix}.
\end{definition}

\begin{lemma}[The primal virtual labels fuse to the identity]%
    \label{thm:peps_quantum_double_primal_fusion}
    \lean{TNLean.PEPS.quantumDoublePrimalTensor_ne_zero_mul_eq_one,
        TNLean.PEPS.quantumDoublePrimalTensor_ne_zero_map_mul_eq_one,
        TNLean.PEPS.quantumDoublePrimalTensor_smul_map_mul,
        TNLean.PEPS.character_mul_quantumDoublePrimalTensor}
    \leanok
    \uses{def:peps_quantum_double_primal}
    On the support (\ref{eq:peps_qd_primal_labels}), $t\,r\,b^{-1}l^{-1}=1$. Hence for every representation $\pi$ of $G$,
    \begin{align}
        \pi(t)\pi(r)\pi(b)^{-1}\pi(l)^{-1} & =I,
        \label{eq:peps_qd_primal_fusion}
    \end{align}
    and for a character $\chi$, $\chi(t)\chi(r)\chi(b)^{-1}\chi(l)^{-1}A^{g_1,g_2,g_3,g_4}_{t,r,b,l}=A^{g_1,g_2,g_3,g_4}_{t,r,b,l}$ \cite[Appendix~A, ``The Toric Code and quantum double models'']{Cirac2021Matrix}. The resulting network string symmetry is Lemma~\ref{thm:peps_quantum_double_primal_network}.
\end{lemma}

\begin{proof}\leanok
    From (\ref{eq:peps_qd_primal_labels}), $t\,r\,b^{-1}l^{-1}=(g_1g_2)(g_2^{-1}g_3)(g_4g_3)^{-1}(g_1g_4^{-1})^{-1}=1$; multiplicativity of $\pi$ gives (\ref{eq:peps_qd_primal_fusion}).
\end{proof}

\begin{lemma}[Virtual symmetry of the primal quantum-double network]%
    \label{thm:peps_quantum_double_primal_network}
    \lean{TNLean.PEPS.quantumDoublePrimalTensor_ne_zero_mul_eq_mul,
        TNLean.PEPS.quantumDoublePrimalSiteTensor,
        TNLean.PEPS.quantumDoublePrimalPEPS,
        TNLean.PEPS.sum_smul_torusWestSouthString_quantumDoublePrimalPEPS_eq}
    \leanok
    \uses{def:peps_quantum_double_primal, lem:peps_torus_string_pull_through}
    Place the primal tensor of Definition~\ref{def:peps_quantum_double_primal} at every site of a $w\times h$ torus ($w,h\ge2$; four distinct bonds per site for $w,h\ge3$ \cite{gap:rmp_peps_examples_small_torus}). For any representation $\pi$ of $G$, inserting $\pi(g)$ along the west and south boundary of a site rectangle yields the same contracted network as inserting $\pi(g)$ along its north and east boundary \cite[Appendix~A, ``The Toric Code and quantum double models'']{Cirac2021Matrix}; around a single site this reproduces (\ref{eq:peps_qd_primal_fusion}).
\end{lemma}

\begin{proof}\leanok
    On the support (\ref{eq:peps_qd_primal_labels}), $tr=lb$, so $\pi(t)\pi(r)=\pi(l)\pi(b)$; apply Lemma~\ref{lem:peps_torus_string_pull_through} with $X_g=\pi(g)$.
\end{proof}
